\section{Detection results per model}
\label{supp:results}

\subsection{Detection results per poison rate}
\label{app:results}

\Cref{tab:results-rate-1,tab:results-rate-5,tab:results-rate-10} give the detection results at 1\%, 5\% and 10\% poisoning, one row per backdoored model, grouped by dataset. ASR is the attack success rate, the share of triggered non-target images classified as the target, and CA is the clean accuracy on the test set. The first block of detection columns is PSBD-TM (token masking at the attention input) and the second is PSBD-RD (dropout on the residual stream after both additions). AUROC is one-sided and never flipped, so a value below chance means the statistic ordered poisoned and clean inputs the wrong way round. TPR@FPR is the share of triggered inputs flagged when the threshold is set on the \HeldoutSize{} clean validation images so that 10\% or 20\% of them are flagged.

A model whose attack success is below \AsrBar{} shows its ASR and CA but no detection numbers, since a detector cannot be credited or blamed for a backdoor that was never planted. BadNets, Blend, BPP and LF implant at every rate on every dataset. TaCT implants everywhere except GTSRB at 1\%. WaNet implants only at 5\% and 10\% and not on CIFAR-100. SIG implants only on CIFAR-10 at 10\%, and Adaptive-Blend never reaches the bar. \Cref{tab:results-benign} gives the clean accuracy of the benign reference model of each dataset, which every clean-accuracy drop in the paper is measured against.

% generated by scripts/paper/tab_results_by_rate.py from results/coverage/coverage.json, results/<folder>/psbd_metrics.json at 2026-09-29T18:10:31.697606+00:00, commit d73d3277d007fd75b3a4166a7911b1e07a2761b7-dirty. Do not edit by hand.
\begin{table*}[htbp]
\centering
\small
\setlength{\tabcolsep}{4pt}
\caption{Detection results at 1\% poisoning on ViT-B/16. A model that is not a successful backdoor (attack success below the bar, or clean accuracy more than 2 points below the benign model) prints -- in every detection column.}
\label{tab:results-rate-1}
\begin{tabular}{lrrrrrrrr}
\toprule
 & & & \multicolumn{3}{c}{PSBD-TM} & \multicolumn{3}{c}{PSBD-RD} \\
\cmidrule(lr){4-6} \cmidrule(lr){7-9}
Attack & ASR & CA & AUROC & \multicolumn{2}{c}{TPR@FPR} & AUROC & \multicolumn{2}{c}{TPR@FPR} \\
\cmidrule(lr){5-6} \cmidrule(lr){8-9}
 & & & & 10\% & 20\% & & 10\% & 20\% \\
\midrule
\multicolumn{9}{l}{\textit{CIFAR-10}} \\
BadNets & 0.997 & 0.952 & \textbf{0.982} & 0.959 & 0.978 & 0.309 & 0.015 & 0.046 \\
Blend & 1.000 & 0.954 & 0.922 & 0.773 & 0.869 & \textbf{0.990} & 0.983 & 0.992 \\
LF & 0.982 & 0.953 & \textbf{0.979} & 0.954 & 0.975 & 0.941 & 0.753 & 0.898 \\
BPP & 0.982 & 0.952 & \textbf{0.986} & 0.956 & 0.976 & 0.970 & 0.897 & 0.952 \\
WaNet & 0.112 & 0.954 & -- & -- & -- & -- & -- & -- \\
TaCT & 0.985 & 0.938 & \textbf{0.979} & 0.005 & 0.009 & 0.464 & 0.013 & 0.038 \\
SIG & 0.340 & 0.951 & -- & -- & -- & -- & -- & -- \\
Adaptive-Blend & 0.622 & 0.953 & -- & -- & -- & -- & -- & -- \\
\midrule
\multicolumn{9}{l}{\textit{CIFAR-100}} \\
BadNets & 1.000 & 0.825 & \textbf{0.988} & 0.999 & 1.000 & 0.744 & 0.244 & 0.476 \\
Blend & 0.989 & 0.829 & \textbf{0.976} & 0.959 & 0.983 & 0.792 & 0.347 & 0.598 \\
LF & 0.942 & 0.830 & 0.956 & 0.942 & 0.947 & \textbf{0.956} & 0.903 & 0.940 \\
BPP & 0.914 & 0.824 & 0.921 & 0.849 & 0.919 & \textbf{0.950} & 0.914 & 0.929 \\
WaNet & 0.057 & 0.822 & -- & -- & -- & -- & -- & -- \\
TaCT & 0.989 & 0.820 & -- & -- & -- & -- & -- & -- \\
SIG & 0.250 & 0.817 & -- & -- & -- & -- & -- & -- \\
Adaptive-Blend & 0.536 & 0.823 & -- & -- & -- & -- & -- & -- \\
\midrule
\multicolumn{9}{l}{\textit{GTSRB}} \\
BadNets & 1.000 & 0.994 & \textbf{0.999} & 0.996 & 0.999 & 0.222 & 0.001 & 0.008 \\
Blend & 1.000 & 0.991 & \textbf{1.000} & 1.000 & 1.000 & 0.999 & 1.000 & 1.000 \\
LF & 0.979 & 0.990 & 0.976 & 0.926 & 0.973 & \textbf{0.987} & 0.977 & 0.985 \\
BPP & 0.868 & 0.986 & \textbf{0.898} & 0.717 & 0.809 & 0.755 & 0.128 & 0.507 \\
WaNet & 0.119 & 0.979 & -- & -- & -- & -- & -- & -- \\
TaCT & 1.000 & 0.994 & pending & pending & pending & pending & pending & pending \\
SIG & 0.356 & 0.981 & -- & -- & -- & -- & -- & -- \\
Label-Consistent & 0.689 & 0.991 & -- & -- & -- & -- & -- & -- \\
Adaptive-Blend & 0.560 & 0.983 & -- & -- & -- & -- & -- & -- \\
\midrule
\multicolumn{9}{l}{\textit{Tiny ImageNet}} \\
BadNets & 0.999 & 0.752 & \textbf{0.980} & 0.989 & 0.997 & 0.892 & 0.700 & 0.836 \\
Blend & 0.999 & 0.760 & 0.992 & 0.985 & 0.992 & \textbf{0.996} & 0.994 & 0.998 \\
LF & 0.922 & 0.762 & 0.941 & 0.929 & 0.935 & \textbf{0.952} & 0.918 & 0.934 \\
BPP & 0.971 & 0.748 & 0.980 & 0.972 & 0.975 & \textbf{0.981} & 0.972 & 0.974 \\
WaNet & 0.178 & 0.747 & -- & -- & -- & -- & -- & -- \\
TaCT & 0.976 & 0.750 & -- & -- & -- & -- & -- & -- \\
SIG & 0.076 & 0.762 & -- & -- & -- & -- & -- & -- \\
Adaptive-Blend & 0.554 & 0.747 & -- & -- & -- & -- & -- & -- \\
\bottomrule
\end{tabular}
\end{table*}

% generated by scripts/paper/tab_results_by_rate.py from results/coverage/coverage.json, results/<folder>/psbd_metrics.json at 2026-09-30T13:15:46.759123+00:00, commit 257fa16f3d6cf6ee17ee984d58eef143db7ecc4c-dirty. Do not edit by hand.
\begin{table*}[htbp]
\centering
\small
\setlength{\tabcolsep}{4pt}
\caption{Detection results at 5\% poisoning on ViT-B/16. A model that is not a successful backdoor (attack success below the bar, or clean accuracy more than 2 points below the benign model) prints -- in every detection column.}
\label{tab:results-rate-5}
\begin{tabular}{lrrrrrrrr}
\toprule
 & & & \multicolumn{3}{c}{PSBD-TM} & \multicolumn{3}{c}{PSBD-RD} \\
\cmidrule(lr){4-6} \cmidrule(lr){7-9}
Attack & ASR & CA & AUROC & \multicolumn{2}{c}{TPR@FPR} & AUROC & \multicolumn{2}{c}{TPR@FPR} \\
\cmidrule(lr){5-6} \cmidrule(lr){8-9}
 & & & & 10\% & 20\% & & 10\% & 20\% \\
\midrule
\multicolumn{9}{l}{\textit{CIFAR-10}} \\
BadNets & 1.000 & 0.945 & \textbf{0.992} & 0.965 & 0.995 & 0.630 & 0.031 & 0.168 \\
Blend & 1.000 & 0.957 & 0.971 & 0.908 & 0.989 & \textbf{0.984} & 0.973 & 0.995 \\
LF & 0.997 & 0.944 & \textbf{0.991} & 0.983 & 0.994 & 0.959 & 0.837 & 0.947 \\
BPP & 0.987 & 0.950 & 0.801 & 0.587 & 0.655 & \textbf{0.880} & 0.673 & 0.788 \\
WaNet & 0.961 & 0.913 & -- & -- & -- & -- & -- & -- \\
TaCT & 0.996 & 0.946 & \textbf{0.966} & 0.000 & 0.031 & 0.834 & 0.006 & 0.020 \\
SIG & 0.599 & 0.953 & -- & -- & -- & -- & -- & -- \\
Adaptive-Blend & 0.594 & 0.946 & -- & -- & -- & -- & -- & -- \\
\midrule
\multicolumn{9}{l}{\textit{CIFAR-100}} \\
BadNets & 1.000 & 0.825 & \textbf{0.994} & 0.999 & 1.000 & 0.826 & 0.335 & 0.672 \\
Blend & 1.000 & 0.830 & 0.979 & 0.941 & 0.976 & \textbf{0.999} & 0.999 & 1.000 \\
LF & 0.953 & 0.829 & \textbf{0.967} & 0.954 & 0.959 & 0.966 & 0.934 & 0.955 \\
BPP & 0.988 & 0.833 & 0.985 & 0.961 & 0.977 & \textbf{0.988} & 0.973 & 0.984 \\
WaNet & 0.643 & 0.809 & -- & -- & -- & -- & -- & -- \\
TaCT & 0.980 & 0.833 & pending & pending & pending & pending & pending & pending \\
SIG & 0.249 & 0.812 & -- & -- & -- & -- & -- & -- \\
Adaptive-Blend & 0.579 & 0.812 & -- & -- & -- & -- & -- & -- \\
\midrule
\multicolumn{9}{l}{\textit{GTSRB}} \\
BadNets & 1.000 & 0.986 & \textbf{1.000} & 1.000 & 1.000 & 0.889 & 0.659 & 0.830 \\
Blend & 1.000 & 0.977 & \textbf{0.999} & 1.000 & 1.000 & 0.899 & 0.714 & 0.834 \\
LF & 0.999 & 0.974 & \textbf{0.998} & 0.999 & 0.999 & 0.997 & 0.997 & 0.999 \\
BPP & 0.991 & 0.978 & 0.992 & 0.978 & 0.987 & \textbf{0.994} & 0.988 & 0.992 \\
WaNet & 0.830 & 0.971 & -- & -- & -- & -- & -- & -- \\
TaCT & 1.000 & 0.986 & \textbf{0.942} & 0.502 & 0.818 & 0.392 & 0.010 & 0.041 \\
SIG & 0.673 & 0.990 & -- & -- & -- & -- & -- & -- \\
Label-Consistent & 0.958 & 0.981 & \textbf{0.969} & 0.950 & 0.963 & 0.444 & 0.031 & 0.088 \\
Adaptive-Blend & 0.777 & 0.981 & -- & -- & -- & -- & -- & -- \\
\midrule
\multicolumn{9}{l}{\textit{Tiny ImageNet}} \\
BadNets & 1.000 & 0.752 & \textbf{0.984} & 0.999 & 1.000 & 0.837 & 0.566 & 0.739 \\
Blend & 0.999 & 0.754 & 0.997 & 0.997 & 0.999 & \textbf{0.999} & 0.999 & 0.999 \\
LF & 0.987 & 0.753 & 0.987 & 0.987 & 0.988 & \textbf{0.991} & 0.987 & 0.989 \\
BPP & 0.985 & 0.746 & 0.987 & 0.984 & 0.986 & \textbf{0.987} & 0.977 & 0.987 \\
WaNet & 0.922 & 0.748 & 0.930 & 0.924 & 0.928 & \textbf{0.953} & 0.922 & 0.938 \\
TaCT & 0.987 & 0.757 & pending & pending & pending & pending & pending & pending \\
SIG & 0.432 & 0.748 & -- & -- & -- & -- & -- & -- \\
Adaptive-Blend & 0.783 & 0.745 & -- & -- & -- & -- & -- & -- \\
\bottomrule
\end{tabular}
\end{table*}

% generated by scripts/paper/tab_results_by_rate.py from results/coverage/coverage.json, results/<folder>/psbd_metrics.json at 2026-09-29T16:31:44.769994+00:00, commit 19488b81358b06040ba96f361d5061b2981f1f25-dirty. Do not edit by hand.
\begin{table*}[htbp]
\centering
\small
\setlength{\tabcolsep}{4pt}
\caption{Detection results at 10\% poisoning on ViT-B/16.}
\label{tab:results-rate-10}
\begin{tabular}{lrrrrrrrr}
\toprule
 & & & \multicolumn{3}{c}{PSBD-TM} & \multicolumn{3}{c}{PSBD-RD} \\
\cmidrule(lr){4-6} \cmidrule(lr){7-9}
Attack & ASR & CA & AUROC & \multicolumn{2}{c}{TPR@FPR} & AUROC & \multicolumn{2}{c}{TPR@FPR} \\
\cmidrule(lr){5-6} \cmidrule(lr){8-9}
 & & & & 10\% & 20\% & & 10\% & 20\% \\
\midrule
\multicolumn{9}{l}{\textit{CIFAR-10}} \\
BadNets & 1.000 & 0.942 & \textbf{0.991} & 0.935 & 1.000 & 0.720 & 0.113 & 0.301 \\
Blend & 1.000 & 0.951 & 0.906 & 0.754 & 0.849 & \textbf{0.991} & 0.980 & 0.992 \\
LF & 0.999 & 0.951 & \textbf{0.993} & 0.995 & 0.999 & 0.977 & 0.908 & 0.972 \\
BPP & 0.994 & 0.954 & 0.871 & 0.736 & 0.830 & \textbf{0.875} & 0.720 & 0.799 \\
WaNet & 0.890 & 0.946 & 0.459 & 0.087 & 0.268 & \textbf{0.947} & 0.846 & 0.914 \\
TaCT & 1.000 & 0.858 & -- & -- & -- & -- & -- & -- \\
SIG & 0.901 & 0.846 & 0.418 & 0.071 & 0.097 & \textbf{0.919} & 0.810 & 0.878 \\
Adaptive-Blend & 0.622 & 0.951 & -- & -- & -- & -- & -- & -- \\
\midrule
\multicolumn{9}{l}{\textit{CIFAR-100}} \\
BadNets & 1.000 & 0.821 & \textbf{0.998} & 1.000 & 1.000 & 0.497 & 0.114 & 0.234 \\
Blend & 1.000 & 0.816 & 0.998 & 1.000 & 1.000 & \textbf{1.000} & 1.000 & 1.000 \\
LF & 0.995 & 0.823 & \textbf{0.995} & 0.995 & 0.995 & 0.982 & 0.959 & 0.983 \\
BPP & 0.991 & 0.817 & 0.985 & 0.963 & 0.975 & \textbf{0.985} & 0.960 & 0.979 \\
WaNet & 0.793 & 0.818 & -- & -- & -- & -- & -- & -- \\
TaCT & 0.989 & 0.821 & -- & -- & -- & -- & -- & -- \\
SIG & 0.178 & 0.820 & -- & -- & -- & -- & -- & -- \\
Adaptive-Blend & 0.604 & 0.808 & -- & -- & -- & -- & -- & -- \\
\midrule
\multicolumn{9}{l}{\textit{GTSRB}} \\
BadNets & 1.000 & 0.972 & \textbf{1.000} & 1.000 & 1.000 & 0.999 & 0.999 & 1.000 \\
Blend & 1.000 & 0.988 & \textbf{1.000} & 1.000 & 1.000 & 0.998 & 0.996 & 0.998 \\
LF & 0.999 & 0.988 & 0.998 & 0.998 & 0.999 & \textbf{0.999} & 0.999 & 0.999 \\
BPP & 0.995 & 0.986 & \textbf{0.989} & 0.971 & 0.986 & 0.977 & 0.971 & 0.992 \\
WaNet & 0.947 & 0.958 & 0.948 & 0.949 & 0.952 & \textbf{0.960} & 0.950 & 0.960 \\
TaCT & 1.000 & 0.931 & -- & -- & -- & -- & -- & -- \\
Adaptive-Blend & 0.838 & 0.976 & -- & -- & -- & -- & -- & -- \\
\midrule
\multicolumn{9}{l}{\textit{Tiny ImageNet}} \\
BadNets & 1.000 & 0.752 & \textbf{0.993} & 1.000 & 1.000 & 0.973 & 0.951 & 0.985 \\
Blend & 1.000 & 0.749 & 0.998 & 0.999 & 0.999 & \textbf{0.999} & 0.999 & 1.000 \\
LF & 0.973 & 0.740 & 0.973 & 0.973 & 0.977 & \textbf{0.982} & 0.970 & 0.976 \\
BPP & 0.996 & 0.757 & 0.984 & 0.981 & 0.991 & \textbf{0.997} & 0.995 & 0.996 \\
WaNet & 0.968 & 0.747 & 0.950 & 0.932 & 0.970 & \textbf{0.966} & 0.924 & 0.963 \\
TaCT & 0.929 & 0.756 & -- & -- & -- & -- & -- & -- \\
SIG & 0.122 & 0.747 & -- & -- & -- & -- & -- & -- \\
Adaptive-Blend & 0.505 & 0.753 & -- & -- & -- & -- & -- & -- \\
\bottomrule
\end{tabular}
\end{table*}

% generated by scripts/paper/tab_results_by_rate.py from results/coverage/coverage.json at 2026-09-24T05:54:57.120883+00:00, commit 872fbbbd1f025c8c9924211876e2cda72fc45aab-dirty. Do not edit by hand.
\begin{table}[htbp]
\centering
\small
\caption{Benign reference models.}
\label{tab:results-benign}
\begin{adjustbox}{max width=\linewidth}
\begin{tabular}{lrr}
\toprule
Dataset & Clean Accuracy & Attacks Above ASR Bar \\
\midrule
CIFAR-10 & 0.953 & 17 \\
CIFAR-100 & 0.811 & 12 \\
GTSRB & 0.991 & 16 \\
Tiny ImageNet & 0.755 & 14 \\
\bottomrule
\end{tabular}
\end{adjustbox}
\end{table}

\FloatBarrier

\subsection{The Label-Consistent epsilon pilot}
\label{supp:lc-epsilon}

\Cref{tab:clean-label-lc-epsilon} gives the pilot behind the epsilon choice of \cref{app:attacks}, with one run per dataset and epsilon at each dataset's highest reachable clean-label rate.

% generated by scripts/paper/tab_clean_label.py from checkpoints/<folder>/args.json, configs/psbd_basis.json at 2026-09-30T14:53:39.015080+00:00, commit 1fd38381a58031e98490ed24679d373b446725fb-dirty. Do not edit by hand.
\begin{table}[htbp]
\centering
\small
\caption{The Label-Consistent epsilon pilot with Turner's adversarial bases, seed 0 at each dataset's highest reachable clean-label rate. A run passes when it clears the attack success bar inside the clean-accuracy drop bar, and the chosen epsilon is the smallest passing one.}
\label{tab:clean-label-lc-epsilon}
\begin{adjustbox}{max width=\linewidth}
\begin{tabular}{lrrrrrl}
\toprule
dataset & rate & epsilon/255 & ASR & CA & dCA & passes \\
\midrule
CIFAR-10 & 0.1 & 8 & 0.970 & 0.859 & $-$0.094 & no \\
CIFAR-10 & 0.1 & 16 & 0.860 & 0.856 & $-$0.097 & no \\
CIFAR-10 & 0.1 & 32 & 0.393 & 0.856 & $-$0.096 & no \\
CIFAR-100 & 0.01 & 8 & 0.700 & 0.824 & +0.014 & no \\
CIFAR-100 & 0.01 & 16 & 0.756 & 0.826 & +0.015 & no \\
CIFAR-100 & 0.01 & 32 & 0.048 & 0.824 & +0.013 & no \\
GTSRB & 0.05 & 8 & 0.437 & 0.093 & $-$0.898 & diverged \\
GTSRB & 0.05 & 16 & 0.937 & 0.989 & $-$0.002 & yes \\
GTSRB & 0.05 & 32 & 0.321 & 0.081 & $-$0.911 & diverged \\
Tiny ImageNet & 0.005 & 8 & 0.614 & 0.747 & $-$0.008 & no \\
Tiny ImageNet & 0.005 & 16 & 0.508 & 0.751 & $-$0.004 & no \\
Tiny ImageNet & 0.005 & 32 & 0.448 & 0.752 & $-$0.003 & no \\
\bottomrule
\end{tabular}
\end{adjustbox}
\end{table}

\FloatBarrier
\clearpage


\section{Every placement on ViT-B/16}
\label{supp:basis}

\subsection{The placements we swept}
\label{app:basis}

\Cref{tab:basis-declaration} lists every placement we swept on every model. A placement pairs a site (where the perturbation is injected, \cref{fig:vit-block}) with a perturbation (what is injected there). It acts in every block unless a band of blocks is named. The family column groups sites by where they sit relative to the residual stream. Input-side sites perturb a sublayer's input before its LayerNorm, and residual-adjacent sites perturb a branch output just before its residual addition or the stream just after it. The structured sites perturb attention heads and MLP hidden units, and the ported sites come from other detectors, scaling input pixels as SCALE-UP does and scaling a LayerNorm's gain as IBD-PSC does. Masks and dropout take a share of the tensor as their rate, and Gaussian noise a multiple of the tensor's standard deviation.

% generated by scripts/paper/app_basis.py from configs/psbd_basis.json at 2026-09-29T16:06:21.381096+00:00, commit 2e59453ac935840c88040509aab36fb7e2c4ac3d-dirty. Do not edit by hand.
\begin{table}[htbp]
\centering
\small
\caption{Every placement we swept, with the blocks it acts in, its family and the number and range of rates in its ladder.}
\label{tab:basis-declaration}
\begin{adjustbox}{max width=\linewidth}
\begin{tabular}{llll}
\toprule
placement & blocks & family & rates \\
\midrule
token mask, attention input & all & input side & 10 (0.05 to 0.9) \\
channel mask, attention input & all & input side & 10 (0.05 to 0.9) \\
noise, attention input & all & input side & 12 (0.05 to 8) \\
token mask, MLP input & all & input side & 10 (0.05 to 0.9) \\
token mask, both sublayer inputs & all & input side & 9 (0.1 to 0.9) \\
scale up, input pixels & all & input side & 11 (0.5 to 25) \\
noise, MLP input after norm & all & input side & 12 (0.05 to 8) \\
token mask, attention output before the add & all & residual adjacent & 10 (0.05 to 0.9) \\
token mask, stream after the attention add & all & residual adjacent & 10 (0.05 to 0.9) \\
dropout, stream after the attention add & all & residual adjacent & 9 (0.1 to 0.9) \\
dropout, before both residual adds & all & residual adjacent & 9 (0.1 to 0.9) \\
dropout, after both residual adds & all & residual adjacent & 16 (0.005 to 0.9) \\
channel mask, MLP neurons & all & structured & 10 (0.05 to 0.9) \\
gain scale, MLP norm output & all & ported & 11 (0.25 to 15) \\
token mask, attention input, blocks 5 to 8 & 5-8 & depth band & 12 (0.05 to 0.99) \\
token mask, attention input, blocks 9 to 12 & 9-12 & depth band & 12 (0.05 to 0.99) \\
dropout, before both residual adds, blocks 1 to 4 & 1-4 & depth band & 11 (0.1 to 0.99) \\
dropout, before both residual adds, blocks 5 to 8 & 5-8 & depth band & 11 (0.1 to 0.99) \\
dropout, before both residual adds, blocks 9 to 12 & 9-12 & depth band & 11 (0.1 to 0.99) \\
dropout, attention input & all & input side & 9 (0.1 to 0.9) \\
dropout, attention input after norm & all & input side & 9 (0.1 to 0.9) \\
dropout, MLP input & all & input side & 9 (0.1 to 0.9) \\
dropout, embedding output & all & input side & 9 (0.1 to 0.9) \\
token mask, attention input, blocks 1 to 4 & 1-4 & depth band & 12 (0.05 to 0.99) \\
dropout, attention output before the add & all & residual adjacent & 9 (0.1 to 0.9) \\
channel mask, attention output before the add & all & residual adjacent & 10 (0.05 to 0.9) \\
noise, attention output before the add & all & residual adjacent & 12 (0.05 to 8) \\
\bottomrule
\end{tabular}
\end{adjustbox}
\end{table}

\FloatBarrier

\subsection{Every placement ranked}
\label{supp:ranking}

\Cref{tab:basis-ranking} ranks the placements by mean AUROC and gives each one's lowest single-model AUROC, and \cref{fig:basis} plots the ranking. The four best placements all mask whole tokens. PSBD-TM and its twin on the attention branch output sit at the top together, and \cref{app:protocol} says why we recommend the attention input.

% generated by scripts/paper/app_basis.py from results/coverage/coverage.json, configs/psbd_basis.json, results/<folder>/psbd_metrics.json (62 cells) at 2026-09-30T12:57:05.191495+00:00, commit 257fa16f3d6cf6ee17ee984d58eef143db7ecc4c-dirty. Do not edit by hand.
\begin{table}[htbp]
\centering
\small
\caption{Every placement we swept, ranked by mean AUROC on ViT-B/16. n is the number of models whose rate ladder reaches the shift ratio target, and the last column is the lowest single-model AUROC.}
\label{tab:basis-ranking}
\begin{adjustbox}{max width=\linewidth}
\begin{tabular}{rllrrr}
\toprule
Rank & Placement & Family & n & AUROC & Lowest AUROC \\
\midrule
1 & token mask, attention input, blocks 9 to 12 & depth band & 25 & 0.955 & 0.802 \\
2 & token mask, attention output before the add & residual adjacent & 56 & 0.952 & 0.537 \\
3 & token mask, attention input & input side & 56 & 0.951 & 0.266 \\
4 & gain scale, MLP norm output & ported & 56 & 0.938 & 0.307 \\
5 & token mask, attention input, blocks 1 to 4 & depth band & 12 & 0.932 & 0.753 \\
6 & channel mask, attention input & input side & 56 & 0.919 & 0.579 \\
7 & token mask, both sublayer inputs & input side & 56 & 0.911 & 0.010 \\
8 & dropout, before both residual adds, blocks 5 to 8 & depth band & 56 & 0.908 & 0.365 \\
9 & noise, MLP input after norm & input side & 56 & 0.906 & 0.253 \\
10 & dropout, attention input & input side & 56 & 0.905 & 0.178 \\
11 & dropout, before both residual adds, blocks 9 to 12 & depth band & 41 & 0.900 & 0.353 \\
12 & dropout, after both residual adds & residual adjacent & 56 & 0.876 & 0.222 \\
13 & dropout, before both residual adds & residual adjacent & 56 & 0.874 & 0.327 \\
14 & token mask, attention input, blocks 5 to 8 & depth band & 37 & 0.872 & 0.116 \\
15 & dropout, stream after the attention add & residual adjacent & 56 & 0.846 & 0.190 \\
16 & dropout, attention input after norm & input side & 56 & 0.846 & 0.165 \\
17 & dropout, embedding output & input side & 56 & 0.846 & 0.423 \\
18 & dropout, attention output before the add & residual adjacent & 56 & 0.840 & 0.299 \\
19 & noise, attention input & input side & 56 & 0.840 & 0.148 \\
20 & noise, attention output before the add & residual adjacent & 56 & 0.834 & 0.055 \\
21 & channel mask, MLP neurons & structured & 56 & 0.826 & 0.317 \\
22 & dropout, before both residual adds, blocks 1 to 4 & depth band & 56 & 0.824 & 0.120 \\
23 & token mask, MLP input & input side & 56 & 0.823 & 0.005 \\
24 & channel mask, attention output before the add & residual adjacent & 56 & 0.810 & 0.153 \\
25 & token mask, stream after the attention add & residual adjacent & 56 & 0.787 & 0.000 \\
26 & dropout, MLP input & input side & 56 & 0.780 & 0.070 \\
27 & scale up, input pixels & input side & 56 & 0.747 & 0.225 \\
\bottomrule
\end{tabular}
\end{adjustbox}
\end{table}


\begin{figure}[htbp]
\centering
\includegraphics{basis_ranking.pdf}
\caption{Mean AUROC of every placement with a complete sweep, read at the rate the original rule assigns.}
\label{fig:basis}
\end{figure}
\FloatBarrier

\subsection{Depth bands at the adaptive rule}
\label{sec:results:depth}

PSBD-TM perturbs all \VitBlocks{} blocks. \Cref{tab:staircase-bands} restricts it to a third of the network at the adaptive rule. Restricting the mask to blocks 5 to 8 costs \StaircaseBandsBeforeAttentionNormBlocksFiveEightTokenMaskGain{} and to blocks 9 to 12 costs \StaircaseBandsBeforeAttentionNormBlocksNineOneTwoTokenMaskGain{}. Few models reach the shift target with the first third alone, so \cref{app:bands} compares the bands at matched disturbance, where the first third loses most.

% generated by scripts/paper/tab_staircase.py from results/coverage/coverage.json, results/<folder>/psbd_metrics.json (54 models) at 2026-09-29T18:12:06.471752+00:00, commit d73d3277d007fd75b3a4166a7911b1e07a2761b7-dirty. Do not edit by hand.
\begin{table*}[htbp]
\centering
\small
\caption{Token masking at the attention input of ViT-B/16, restricted to bands of blocks.}
\label{tab:staircase-bands}
\begin{tabular}{lrrrrl}
\toprule
 & & & \multicolumn{2}{c}{TPR@FPR} & \\
\cmidrule(lr){4-5}
Placement & n & AUROC & 10\% & 20\% & Gain [95\% CI] \\
\midrule
token mask, attention input, all 12 blocks & 54 & 0.963 & 0.887 & 0.916 & -- \\
token mask, attention input, blocks 1 to 4 & 11 & 0.949 & 0.867 & 0.949 & $-$0.031 [$-$0.044, $-$0.019] \\
token mask, attention input, blocks 5 to 8 & 35 & 0.891 & 0.788 & 0.864 & $-$0.087 [$-$0.143, $-$0.038] \\
token mask, attention input, blocks 9 to 12 & 25 & 0.955 & 0.891 & 0.936 & $-$0.023 [$-$0.045, $-$0.004] \\
\bottomrule
\end{tabular}
\end{table*}

\FloatBarrier

\subsection{Whether placements agree on which inputs are fragile}
\label{app:agreement}

If every placement read the same per-sample quantity, a union of placements could add nothing. \Cref{tab:mech-operator-agreement} gives the Spearman correlation between each placement's per-sample PSU and that of PSBD-TM, each read at its own matched rate. \Cref{tab:mech-operator-agreement-family} averages it within and across families of sites.

% generated by scripts/paper/mech_operator_agreement.py from results/<folder>/psbd/<placement>/rate_*_{clean,backdoor}.pt at 2026-09-30T13:04:41.000399+00:00, commit 257fa16f3d6cf6ee17ee984d58eef143db7ecc4c-dirty. Do not edit by hand.
\begin{table}[htbp]
\centering
\small
\caption{Whether placements agree on which inputs are fragile. Mean Spearman correlation, over the models whose attack succeeded, between each placement's per-sample fractional PSU and that of token masking at the attention input, each at its own matched 0.6 rate. n is the number of models carrying both placements at a usable rate.}
\label{tab:mech-operator-agreement}
\begin{adjustbox}{max width=\linewidth}
\begin{tabular}{lrrrrr}
\toprule
placement & family & rho, clean & n & rho, triggered & n \\
\midrule
channel mask, attention input & input side & 0.633 & 56 & 0.252 & 56 \\
noise, attention input & input side & 0.384 & 56 & 0.235 & 56 \\
token mask, MLP input & input side & 0.564 & 56 & 0.442 & 56 \\
token mask, both sublayer inputs & input side & 0.812 & 56 & 0.577 & 56 \\
scale up, input pixels & input side & 0.213 & 56 & 0.097 & 56 \\
noise, MLP input after norm & input side & 0.545 & 56 & 0.248 & 56 \\
token mask, attention output before the add & residual adjacent & 0.538 & 56 & 0.416 & 56 \\
token mask, stream after the attention add & residual adjacent & 0.577 & 56 & 0.443 & 56 \\
dropout, stream after the attention add & residual adjacent & 0.519 & 56 & 0.277 & 56 \\
dropout, before both residual adds & residual adjacent & 0.598 & 56 & 0.275 & 56 \\
dropout, after both residual adds & residual adjacent & 0.613 & 56 & 0.315 & 56 \\
channel mask, MLP neurons & structured & 0.504 & 56 & 0.173 & 56 \\
gain scale, MLP norm output & ported & 0.372 & 56 & 0.316 & 56 \\
token mask, attention input, blocks 5 to 8 & depth band & 0.570 & 56 & 0.434 & 56 \\
token mask, attention input, blocks 9 to 12 & depth band & 0.686 & 56 & 0.126 & 56 \\
dropout, before both residual adds, blocks 1 to 4 & depth band & 0.560 & 56 & 0.234 & 56 \\
dropout, before both residual adds, blocks 5 to 8 & depth band & 0.650 & 56 & 0.271 & 56 \\
dropout, before both residual adds, blocks 9 to 12 & depth band & 0.439 & 56 & 0.217 & 56 \\
dropout, attention input & input side & 0.660 & 56 & 0.272 & 56 \\
dropout, attention input after norm & input side & 0.407 & 56 & 0.180 & 56 \\
dropout, MLP input & input side & 0.479 & 56 & 0.186 & 56 \\
dropout, embedding output & input side & 0.391 & 56 & 0.190 & 56 \\
token mask, attention input, blocks 1 to 4 & depth band & 0.519 & 56 & 0.333 & 56 \\
dropout, attention output before the add & residual adjacent & 0.542 & 56 & 0.180 & 56 \\
channel mask, attention output before the add & residual adjacent & 0.514 & 56 & 0.185 & 56 \\
noise, attention output before the add & residual adjacent & 0.566 & 56 & 0.277 & 56 \\
\bottomrule
\end{tabular}
\end{adjustbox}
\end{table}

% generated by scripts/paper/mech_operator_agreement.py from results/<folder>/psbd/<placement>/rate_*_{clean,backdoor}.pt at 2026-09-24T04:33:16.351447+00:00, commit 52e4ba1a7ac941365bc37d68e1e902dad4fb4f74-dirty. Do not edit by hand.
\begin{table}[htbp]
\centering
\small
\caption{Mean per-sample PSU agreement, Spearman correlation averaged over placement pairs and over models, within and across the input-side and residual-adjacent families.}
\label{tab:mech-operator-agreement-family}
\begin{adjustbox}{max width=\linewidth}
\begin{tabular}{lrrrr}
\toprule
comparison & clean rho & n pairs & backdoor rho & n pairs \\
\midrule
within the input-side family & 0.469 & 55 & 0.256 & 55 \\
within the residual-adjacent family & 0.583 & 28 & 0.348 & 28 \\
across the 2 families & 0.512 & 88 & 0.298 & 88 \\
\bottomrule
\end{tabular}
\end{adjustbox}
\end{table}

\FloatBarrier

\subsection{PSU distributions}
\label{app:distributions}

\Cref{fig:psu} shows PSU on six models. In the top row, a BadNets, a BPP and a Blend model, the triggered inputs keep their prediction under the mask and pile up near 0, while the clean inputs lose theirs and pile up near 1. The threshold at the \HeadlineQuantilePercent{} percentile of the clean validation set sits between the two groups. The bottom row shows the three weakest readings. On WaNet and SIG on CIFAR-10 the triggered inputs move as much as the clean ones, so the histograms overlap. On TaCT on Tiny ImageNet both groups sit near 1 and only a thin tail of triggered inputs falls below the threshold.

\begin{figure}[htbp]
\centering
\includegraphics[width=\textwidth]{fig_psu_histograms.pdf}
\caption{PSU under PSBD-TM for clean and triggered inputs, with the threshold at the \HeadlineQuantilePercent{} percentile of the clean validation set (dashed). Top, three models where detection succeeds. Bottom, the three weakest readings.}
\label{fig:psu}
\end{figure}
\FloatBarrier
\clearpage


\section{Every detector per attack}
\label{supp:detectors}

\subsection{AUROC and true-positive rates}
\label{app:detectors}

\Cref{tab:detectors-auroc,tab:detectors-tpr10,tab:detectors-tpr20} and the tables that follow each of them give AUROC and the true-positive rate at 10\% and 20\% false-positive budgets for every defense on every model, one table per dataset. Every detector is read on the same models, splits and thresholds as PSBD. The threshold is set on the \HeldoutSize{} clean validation images so that the named share of them is flagged.

Two ports behave in ways specific to ViTs. SentiNet's Grad-CAM \cite{selvaraju2017gradcam} mask covers 0 to 2\% of the trigger at every block, so its transplant carries almost nothing of the trigger. That would explain a reading at chance, but its mean AUROC is \DetectorsAurocSentinet{}, below chance, so the port carries some signal with the wrong sign, which we report without an explanation. IBD-PSC at the paper's fixed amplification factor of \DetectorsIbdPscFactor{} saturates on a ViT because it keeps every prediction. It reads \DetectorsAurocIbdPsc{}, while the calibrated port that searches the factor per model reads \DetectorsAurocBestCompetitor{}. We therefore count only the calibrated port as the competitor.

% generated by scripts/paper/tab_detectors.py from results/coverage/coverage.json (57 of 59 clearing cells carry every defense), results/*/detectors/*_metrics.json, results/*/psbd_metrics.json at 2026-09-29T16:26:29.298213+00:00, commit 19488b81358b06040ba96f361d5061b2981f1f25-dirty. Do not edit by hand.
% generated by scripts/paper/tab_detectors.py from results/coverage/coverage.json (65 of 65 clearing cells carry every defense), results/*/detectors/*_metrics.json, results/*/psbd_metrics.json at 2026-09-24T04:50:59.131433+00:00, commit 9300709d1f6e7a0ba702f153d77765353ffa1368-dirty. Do not edit by hand.
\begin{table}[htbp]
\centering
\small
\caption{AUROC of every defense on CIFAR-10. Rows are attacks at their poison rate in percent and columns defenses, PSBD-TM first. Best in each row is bold and second best underlined, gray marks a reading below chance and shaded rows are means. PSBD-RD is our adaptation of the original PSBD site to a ViT block. Scale-Up$^\dagger$ is the data-limited variant and IBD-PSC$^\ast$ the calibrated one.}
\label{tab:detectors-auroc}
\begin{adjustbox}{max width=\linewidth}
\begin{tabular}{lrrrrrrrrrrrrr}
\toprule
attack, rate \% & \rotatebox{90}{PSBD-TM} & \rotatebox{90}{PSBD-RD} & \rotatebox{90}{Conf} & \rotatebox{90}{STRIP} & \rotatebox{90}{Scale-Up} & \rotatebox{90}{Scale-Up$^\dagger$} & \rotatebox{90}{IBD-PSC} & \rotatebox{90}{IBD-PSC$^\ast$} & \rotatebox{90}{TeCo} & \rotatebox{90}{CD-L} & \rotatebox{90}{Beatrix} & \rotatebox{90}{TED} & \rotatebox{90}{SentiNet} \\
\midrule
BadNets 1 & \textbf{0.982} & \textcolor{black!45}{0.309} & \textcolor{black!45}{0.479} & 0.906 & \underline{0.977} & 0.909 & 0.558 & 0.975 & 0.891 & 0.958 & 0.969 & 0.814 & 0.644 \\
BadNets 5 & \textbf{0.992} & 0.630 & 0.870 & 0.955 & 0.985 & 0.914 & 0.715 & 0.926 & 0.981 & \underline{0.991} & 0.986 & 0.813 & \textcolor{black!45}{0.433} \\
BadNets 10 & 0.991 & 0.720 & 0.878 & \underline{0.997} & 0.989 & 0.689 & 0.957 & 0.979 & 0.946 & \textbf{1.000} & 0.996 & 0.688 & 0.556 \\
Blend 1 & 0.922 & 0.990 & 0.706 & 0.984 & 0.931 & 0.668 & 0.864 & 0.989 & 0.908 & \underline{0.999} & \textbf{1.000} & 0.830 & \textcolor{black!45}{0.396} \\
Blend 5 & 0.971 & 0.984 & 0.890 & 0.955 & 0.930 & 0.617 & 0.885 & 0.990 & 0.895 & \underline{0.999} & \textbf{1.000} & 0.977 & 0.515 \\
Blend 10 & 0.906 & 0.991 & 0.701 & 0.960 & 0.846 & 0.587 & 0.808 & \underline{0.991} & 0.906 & 0.936 & \textbf{1.000} & 0.964 & \textcolor{black!45}{0.391} \\
BPP 1 & \underline{0.986} & 0.970 & \textcolor{black!45}{0.388} & 0.928 & \textcolor{black!45}{0.389} & \textcolor{black!45}{0.210} & \textcolor{black!45}{0.416} & 0.975 & \textcolor{black!45}{0.447} & 0.951 & \textbf{0.994} & 0.866 & \textcolor{black!45}{0.254} \\
BPP 5 & 0.801 & 0.880 & 0.523 & 0.907 & \textcolor{black!45}{0.413} & \textcolor{black!45}{0.158} & 0.541 & \underline{0.974} & \textcolor{black!45}{0.358} & 0.931 & \textbf{0.998} & 0.874 & \textcolor{black!45}{0.122} \\
BPP 10 & 0.871 & 0.875 & 0.738 & 0.950 & \textcolor{black!45}{0.455} & \textcolor{black!45}{0.221} & 0.841 & \underline{0.991} & \textcolor{black!45}{0.482} & 0.974 & \textbf{0.999} & 0.905 & \textcolor{black!45}{0.273} \\
LF 1 & \textbf{0.979} & \underline{0.941} & \textcolor{black!45}{0.268} & 0.729 & 0.787 & 0.646 & \textcolor{black!45}{0.371} & 0.877 & \textcolor{black!45}{0.467} & 0.806 & 0.940 & 0.808 & \textcolor{black!45}{0.268} \\
LF 5 & \textbf{0.991} & 0.959 & \textcolor{black!45}{0.392} & 0.952 & 0.949 & 0.696 & 0.501 & 0.955 & \textcolor{black!45}{0.285} & 0.923 & \underline{0.986} & 0.775 & \textcolor{black!45}{0.295} \\
LF 10 & \textbf{0.993} & \underline{0.977} & 0.780 & 0.914 & 0.909 & 0.685 & 0.613 & 0.883 & \textcolor{black!45}{0.377} & 0.907 & 0.961 & 0.860 & \textcolor{black!45}{0.284} \\
SIG 10 & \textcolor{black!45}{0.418} & 0.919 & 0.866 & 0.806 & 0.767 & 0.750 & 0.531 & 0.599 & 0.797 & 0.637 & \underline{0.992} & \textbf{0.997} & \textcolor{black!45}{0.005} \\
TaCT 1 & \underline{0.979} & \textcolor{black!45}{0.464} & \textcolor{black!45}{0.175} & 0.509 & 0.653 & \textcolor{black!45}{0.181} & \textcolor{black!45}{0.210} & 0.713 & 0.599 & \textcolor{black!45}{0.374} & \textbf{0.997} & 0.926 & \textcolor{black!45}{0.025} \\
TaCT 5 & 0.966 & 0.834 & 0.728 & 0.596 & \underline{0.981} & 0.876 & 0.850 & 0.855 & 0.881 & 0.710 & \textbf{0.998} & 0.970 & \textcolor{black!45}{0.076} \\
TaCT 10 & 0.963 & \textcolor{black!45}{0.261} & 0.975 & 0.797 & \underline{0.991} & 0.940 & 0.974 & 0.928 & 0.601 & 0.971 & \textbf{0.996} & 0.893 & \textcolor{black!45}{0.011} \\
WaNet 5 & 0.937 & \textbf{0.956} & \textcolor{black!45}{0.480} & 0.517 & 0.929 & \textcolor{black!45}{0.404} & 0.556 & 0.693 & \underline{0.946} & 0.729 & 0.714 & 0.810 & \textcolor{black!45}{0.461} \\
WaNet 10 & \textcolor{black!45}{0.459} & \textbf{0.947} & \textcolor{black!45}{0.376} & 0.504 & 0.567 & \textcolor{black!45}{0.380} & \textcolor{black!45}{0.448} & 0.687 & 0.732 & 0.579 & 0.827 & \underline{0.848} & \textcolor{black!45}{0.353} \\
\rowcolor{black!8}\emph{CIFAR-10 mean} & \underline{0.895} & 0.811 & 0.623 & 0.826 & 0.803 & 0.585 & 0.647 & 0.888 & 0.694 & 0.854 & \textbf{0.964} & 0.868 & \textcolor{black!45}{0.298} \\
\bottomrule
\end{tabular}
\end{adjustbox}
\end{table}

% generated by scripts/paper/tab_detectors.py from results/coverage/coverage.json (65 of 65 clearing cells carry every defense), results/*/detectors/*_metrics.json, results/*/psbd_metrics.json at 2026-09-24T04:51:05.458444+00:00, commit 9300709d1f6e7a0ba702f153d77765353ffa1368-dirty. Do not edit by hand.
\begin{table}[htbp]
\centering
\small
\caption{AUROC of every defense on CIFAR-100. Layout as in \cref{tab:detectors-auroc}.}
\label{tab:detectors-auroc-cifar-100}
\begin{adjustbox}{max width=\linewidth}
\begin{tabular}{lrrrrrrrrrrrrr}
\toprule
attack, rate \% & \rotatebox{90}{PSBD-TM} & \rotatebox{90}{PSBD-RD} & \rotatebox{90}{Conf} & \rotatebox{90}{STRIP} & \rotatebox{90}{Scale-Up} & \rotatebox{90}{Scale-Up$^\dagger$} & \rotatebox{90}{IBD-PSC} & \rotatebox{90}{IBD-PSC$^\ast$} & \rotatebox{90}{TeCo} & \rotatebox{90}{CD-L} & \rotatebox{90}{Beatrix} & \rotatebox{90}{TED} & \rotatebox{90}{SentiNet} \\
\midrule
BadNets 1 & 0.988 & 0.744 & 0.752 & 0.996 & 0.986 & 0.997 & 0.803 & 0.981 & 0.991 & \textbf{1.000} & \underline{0.997} & 0.990 & 0.978 \\
BadNets 5 & 0.994 & 0.826 & 0.980 & 0.993 & 0.989 & 0.992 & 0.959 & 0.972 & 0.994 & \textbf{1.000} & 0.878 & \underline{0.994} & 0.913 \\
BadNets 10 & \underline{0.998} & \textcolor{black!45}{0.497} & 0.965 & 0.995 & 0.992 & 0.998 & 0.982 & 0.990 & 0.993 & \textbf{0.999} & 0.998 & 0.979 & 0.967 \\
Blend 1 & 0.976 & 0.792 & 0.598 & 0.655 & 0.679 & 0.657 & 0.810 & 0.973 & \underline{0.987} & 0.860 & \textbf{0.990} & 0.983 & \textcolor{black!45}{0.317} \\
Blend 5 & 0.979 & \underline{0.999} & 0.958 & 0.897 & 0.782 & 0.730 & 0.931 & 0.987 & 0.976 & 0.953 & 0.984 & \textbf{1.000} & 0.899 \\
Blend 10 & \underline{0.998} & \textbf{1.000} & 0.966 & 0.995 & 0.916 & 0.778 & 0.992 & 0.997 & 0.872 & 0.938 & 0.972 & 0.951 & 0.764 \\
BPP 1 & 0.921 & \underline{0.950} & \textcolor{black!45}{0.375} & 0.819 & 0.530 & \textcolor{black!45}{0.258} & 0.622 & 0.875 & 0.561 & \textbf{0.967} & 0.902 & 0.940 & 0.855 \\
BPP 5 & 0.985 & \underline{0.988} & 0.910 & 0.952 & 0.675 & \textcolor{black!45}{0.244} & 0.870 & 0.937 & 0.917 & \textbf{0.993} & 0.983 & 0.982 & 0.805 \\
BPP 10 & 0.985 & \underline{0.985} & 0.853 & 0.942 & 0.646 & \textcolor{black!45}{0.328} & 0.848 & 0.967 & 0.724 & 0.984 & \textbf{0.986} & 0.981 & 0.602 \\
LF 1 & 0.956 & 0.956 & 0.657 & 0.753 & 0.701 & 0.595 & 0.815 & 0.920 & 0.674 & \textcolor{black!45}{0.383} & \textbf{0.958} & \underline{0.956} & \textcolor{black!45}{0.499} \\
LF 5 & \underline{0.967} & 0.966 & 0.807 & 0.632 & 0.608 & 0.562 & \textcolor{black!45}{0.463} & 0.760 & 0.597 & \textcolor{black!45}{0.261} & 0.957 & \textbf{0.975} & 0.815 \\
LF 10 & \textbf{0.995} & 0.982 & 0.932 & 0.882 & 0.599 & 0.549 & 0.907 & 0.952 & \textcolor{black!45}{0.408} & \textcolor{black!45}{0.377} & 0.972 & \underline{0.983} & 0.730 \\
TaCT 1 & 0.909 & \textcolor{black!45}{0.493} & 0.919 & 0.673 & 0.668 & 0.622 & \textbf{0.935} & \underline{0.931} & 0.908 & 0.772 & 0.604 & 0.650 & \textcolor{black!45}{0.338} \\
TaCT 5 & 0.933 & 0.554 & \textbf{0.998} & 0.975 & \underline{0.997} & 0.995 & 0.983 & 0.980 & 0.921 & 0.956 & 0.658 & 0.636 & \textcolor{black!45}{0.307} \\
TaCT 10 & 0.891 & \textcolor{black!45}{0.488} & \textbf{0.982} & 0.738 & 0.977 & \underline{0.982} & 0.971 & 0.966 & 0.918 & 0.961 & 0.629 & 0.664 & \textcolor{black!45}{0.308} \\
\rowcolor{black!8}\emph{CIFAR-100 mean} & \textbf{0.965} & 0.815 & 0.844 & 0.860 & 0.783 & 0.686 & 0.859 & \underline{0.946} & 0.829 & 0.827 & 0.898 & 0.911 & 0.673 \\
\bottomrule
\end{tabular}
\end{adjustbox}
\end{table}

% generated by scripts/paper/tab_detectors.py from results/coverage/coverage.json (65 of 65 clearing cells carry every defense), results/*/detectors/*_metrics.json, results/*/psbd_metrics.json at 2026-09-24T04:51:11.873676+00:00, commit 9300709d1f6e7a0ba702f153d77765353ffa1368-dirty. Do not edit by hand.
\begin{table}[htbp]
\centering
\small
\caption{AUROC of every defense on GTSRB. Layout as in \cref{tab:detectors-auroc}.}
\label{tab:detectors-auroc-gtsrb}
\begin{adjustbox}{max width=\linewidth}
\begin{tabular}{lrrrrrrrrrrrrr}
\toprule
attack, rate \% & \rotatebox{90}{PSBD-TM} & \rotatebox{90}{PSBD-RD} & \rotatebox{90}{Conf} & \rotatebox{90}{STRIP} & \rotatebox{90}{Scale-Up} & \rotatebox{90}{Scale-Up$^\dagger$} & \rotatebox{90}{IBD-PSC} & \rotatebox{90}{IBD-PSC$^\ast$} & \rotatebox{90}{TeCo} & \rotatebox{90}{CD-L} & \rotatebox{90}{Beatrix} & \rotatebox{90}{TED} & \rotatebox{90}{SentiNet} \\
\midrule
BadNets 1 & 0.999 & \textcolor{black!45}{0.222} & \textcolor{black!45}{0.292} & 0.995 & 0.908 & \textbf{1.000} & 0.630 & 0.929 & 0.982 & 0.980 & \underline{1.000} & 0.999 & \textcolor{black!45}{0.122} \\
BadNets 5 & 1.000 & 0.889 & 0.726 & \textbf{1.000} & 0.914 & \textbf{1.000} & 0.506 & 0.999 & 0.996 & \underline{1.000} & 1.000 & 0.998 & \textcolor{black!45}{0.090} \\
BadNets 10 & \underline{1.000} & 0.999 & \textcolor{black!45}{0.299} & 0.979 & 0.901 & \textbf{1.000} & \textcolor{black!45}{0.328} & 0.993 & 0.954 & 1.000 & 1.000 & 0.999 & \textcolor{black!45}{0.137} \\
Blend 1 & \underline{1.000} & 0.999 & \textcolor{black!45}{0.165} & 0.963 & 0.679 & 0.873 & \textcolor{black!45}{0.410} & 1.000 & 0.771 & 0.936 & \textbf{1.000} & 1.000 & \textcolor{black!45}{0.220} \\
Blend 5 & \underline{0.999} & 0.899 & 0.876 & 0.977 & 0.686 & 0.941 & 0.970 & 0.993 & 0.918 & 0.982 & \textbf{1.000} & 0.996 & \textcolor{black!45}{0.323} \\
Blend 10 & \textbf{1.000} & 0.998 & \textcolor{black!45}{0.413} & 0.966 & 0.692 & 0.944 & 0.993 & \underline{1.000} & 0.831 & 0.998 & 1.000 & 0.999 & \textcolor{black!45}{0.302} \\
BPP 1 & 0.898 & 0.755 & \textcolor{black!45}{0.035} & \textcolor{black!45}{0.148} & \textcolor{black!45}{0.245} & \textcolor{black!45}{0.428} & \textcolor{black!45}{0.071} & 0.921 & \textcolor{black!45}{0.498} & 0.935 & \textbf{0.989} & \underline{0.940} & 0.699 \\
BPP 5 & 0.992 & 0.994 & 0.637 & 0.914 & \textcolor{black!45}{0.144} & \textcolor{black!45}{0.359} & \textcolor{black!45}{0.352} & 0.821 & \textcolor{black!45}{0.284} & 0.985 & \textbf{0.999} & \underline{0.995} & \textcolor{black!45}{0.197} \\
BPP 10 & 0.989 & 0.977 & 0.579 & 0.892 & \textcolor{black!45}{0.148} & \textcolor{black!45}{0.393} & \textcolor{black!45}{0.330} & 0.738 & \textcolor{black!45}{0.350} & 0.967 & \textbf{0.999} & \underline{0.997} & \textcolor{black!45}{0.121} \\
LF 1 & 0.976 & 0.987 & \textcolor{black!45}{0.133} & 0.813 & 0.606 & 0.894 & \textcolor{black!45}{0.226} & 0.954 & 0.661 & \textcolor{black!45}{0.295} & \textbf{0.995} & \underline{0.988} & \textcolor{black!45}{0.126} \\
LF 5 & \underline{0.998} & 0.997 & \textcolor{black!45}{0.491} & 0.942 & 0.742 & 0.917 & \textcolor{black!45}{0.384} & 0.969 & \textcolor{black!45}{0.122} & 0.688 & \textbf{0.999} & 0.993 & \textcolor{black!45}{0.007} \\
LF 10 & 0.998 & 0.999 & 0.598 & 0.960 & 0.656 & 0.870 & 0.705 & 0.951 & \textcolor{black!45}{0.146} & 0.625 & \underline{0.999} & \textbf{1.000} & \textcolor{black!45}{0.473} \\
TaCT 5 & 0.942 & \textcolor{black!45}{0.392} & 0.999 & 0.862 & 0.992 & \underline{1.000} & 0.991 & 0.981 & 0.995 & 0.521 & \textbf{1.000} & 0.999 & \textcolor{black!45}{0.001} \\
TaCT 10 & 0.748 & \textcolor{black!45}{0.295} & 0.999 & 0.979 & 0.994 & \textbf{1.000} & \underline{0.999} & 0.991 & 0.969 & 0.940 & 0.734 & \textcolor{black!45}{0.370} & 0.711 \\
WaNet 10 & 0.948 & \underline{0.960} & \textcolor{black!45}{0.482} & 0.504 & 0.529 & \textcolor{black!45}{0.427} & \textcolor{black!45}{0.475} & 0.643 & 0.948 & 0.549 & 0.667 & \textbf{0.964} & \textcolor{black!45}{0.239} \\
\rowcolor{black!8}\emph{GTSRB mean} & \textbf{0.966} & 0.824 & 0.515 & 0.860 & 0.656 & 0.803 & 0.558 & 0.925 & 0.695 & 0.827 & \underline{0.959} & 0.949 & \textcolor{black!45}{0.251} \\
\bottomrule
\end{tabular}
\end{adjustbox}
\end{table}

% generated by scripts/paper/tab_detectors.py from results/coverage/coverage.json (65 of 65 clearing cells carry every defense), results/*/detectors/*_metrics.json, results/*/psbd_metrics.json at 2026-09-24T04:51:18.433986+00:00, commit 9300709d1f6e7a0ba702f153d77765353ffa1368-dirty. Do not edit by hand.
\begin{table}[htbp]
\centering
\small
\caption{AUROC of every defense on Tiny ImageNet. Layout as in \cref{tab:detectors-auroc}.}
\label{tab:detectors-auroc-tiny-imagenet}
\begin{adjustbox}{max width=\linewidth}
\begin{tabular}{lrrrrrrrrrrrrr}
\toprule
attack, rate \% & \rotatebox{90}{PSBD-TM} & \rotatebox{90}{PSBD-RD} & \rotatebox{90}{Conf} & \rotatebox{90}{STRIP} & \rotatebox{90}{Scale-Up} & \rotatebox{90}{Scale-Up$^\dagger$} & \rotatebox{90}{IBD-PSC} & \rotatebox{90}{IBD-PSC$^\ast$} & \rotatebox{90}{TeCo} & \rotatebox{90}{CD-L} & \rotatebox{90}{Beatrix} & \rotatebox{90}{TED} & \rotatebox{90}{SentiNet} \\
\midrule
BadNets 1 & 0.980 & 0.892 & 0.741 & 0.922 & \underline{0.991} & 0.989 & 0.901 & 0.961 & \textbf{0.993} & 0.708 & 0.851 & 0.930 & \textcolor{black!45}{0.438} \\
BadNets 5 & 0.984 & 0.837 & 0.939 & 0.949 & \underline{0.993} & 0.958 & 0.932 & 0.942 & \textbf{0.995} & 0.839 & 0.763 & 0.863 & \textcolor{black!45}{0.394} \\
BadNets 10 & \underline{0.993} & 0.973 & 0.978 & 0.931 & \textbf{0.995} & 0.940 & 0.949 & 0.957 & 0.991 & 0.846 & 0.557 & 0.840 & \textcolor{black!45}{0.312} \\
Blend 1 & 0.992 & \underline{0.996} & 0.886 & 0.832 & 0.851 & 0.686 & 0.942 & \textbf{0.998} & 0.928 & 0.577 & 0.888 & 0.691 & \textcolor{black!45}{0.370} \\
Blend 5 & 0.997 & \textbf{0.999} & 0.943 & 0.910 & 0.892 & 0.687 & 0.972 & \underline{0.999} & 0.931 & 0.619 & 0.624 & 0.684 & \textcolor{black!45}{0.175} \\
Blend 10 & 0.998 & \textbf{0.999} & 0.995 & 0.975 & 0.950 & 0.724 & 0.993 & \underline{0.999} & 0.955 & \textcolor{black!45}{0.408} & 0.579 & 0.677 & \textcolor{black!45}{0.497} \\
BPP 1 & \underline{0.980} & \textbf{0.981} & 0.833 & 0.953 & \textcolor{black!45}{0.459} & \textcolor{black!45}{0.184} & 0.959 & 0.977 & 0.755 & 0.722 & 0.727 & 0.709 & 0.653 \\
BPP 5 & \underline{0.987} & \textbf{0.987} & 0.909 & 0.940 & \textcolor{black!45}{0.337} & \textcolor{black!45}{0.077} & 0.970 & 0.977 & 0.634 & 0.817 & 0.736 & 0.852 & 0.626 \\
BPP 10 & \underline{0.984} & \textbf{0.997} & 0.858 & 0.964 & \textcolor{black!45}{0.346} & \textcolor{black!45}{0.058} & 0.886 & 0.886 & 0.732 & 0.827 & 0.804 & 0.647 & \textcolor{black!45}{0.417} \\
LF 1 & \underline{0.941} & \textbf{0.952} & 0.767 & 0.608 & 0.710 & 0.628 & 0.876 & 0.876 & 0.692 & 0.783 & 0.632 & 0.710 & \textcolor{black!45}{0.298} \\
LF 5 & \underline{0.987} & \textbf{0.991} & 0.951 & 0.845 & 0.813 & \textcolor{black!45}{0.455} & 0.972 & 0.972 & 0.542 & \textcolor{black!45}{0.483} & 0.507 & 0.557 & \textcolor{black!45}{0.380} \\
LF 10 & \underline{0.973} & \textbf{0.982} & 0.926 & 0.733 & 0.734 & 0.630 & 0.954 & 0.954 & 0.748 & 0.660 & 0.856 & 0.686 & \textcolor{black!45}{0.466} \\
TaCT 1 & 0.575 & \textcolor{black!45}{0.487} & \textbf{0.827} & 0.545 & 0.536 & \textcolor{black!45}{0.429} & 0.773 & \underline{0.776} & 0.618 & \textcolor{black!45}{0.329} & \textcolor{black!45}{0.449} & 0.532 & \textcolor{black!45}{0.411} \\
TaCT 5 & 0.663 & \textcolor{black!45}{0.488} & \underline{0.845} & 0.561 & 0.518 & \textcolor{black!45}{0.496} & \textbf{0.850} & \textbf{0.850} & 0.673 & \textcolor{black!45}{0.252} & 0.645 & \textcolor{black!45}{0.481} & 0.554 \\
TaCT 10 & \underline{0.810} & \textcolor{black!45}{0.417} & \textbf{0.840} & 0.501 & 0.549 & \textcolor{black!45}{0.477} & 0.668 & 0.668 & 0.595 & \textcolor{black!45}{0.364} & \textcolor{black!45}{0.461} & 0.541 & \textcolor{black!45}{0.488} \\
WaNet 5 & \underline{0.930} & \textbf{0.953} & 0.721 & 0.538 & 0.500 & \textcolor{black!45}{0.312} & 0.859 & 0.859 & 0.927 & 0.755 & \textcolor{black!45}{0.228} & 0.638 & \textcolor{black!45}{0.325} \\
WaNet 10 & 0.950 & \textbf{0.966} & 0.836 & 0.532 & 0.549 & \textcolor{black!45}{0.290} & 0.864 & 0.864 & \underline{0.955} & 0.906 & \textcolor{black!45}{0.414} & 0.772 & \textcolor{black!45}{0.364} \\
\rowcolor{black!8}\emph{Tiny ImageNet mean} & \textbf{0.925} & 0.876 & 0.870 & 0.779 & 0.690 & 0.531 & 0.901 & \underline{0.913} & 0.804 & 0.641 & 0.631 & 0.695 & \textcolor{black!45}{0.422} \\
\rowcolor{black!8}\emph{all models} & \textbf{0.935} & 0.832 & 0.714 & 0.829 & 0.735 & 0.644 & 0.742 & \underline{0.916} & 0.754 & 0.786 & 0.860 & 0.851 & \textcolor{black!45}{0.406} \\
\bottomrule
\end{tabular}
\end{adjustbox}
\end{table}


\FloatBarrier
% generated by scripts/paper/tab_detectors.py from results/coverage/coverage.json (56 of 62 successful cells carry every defense), results/*/detectors/*_metrics.json, results/*/psbd_metrics.json at 2026-09-30T13:11:35.728949+00:00, commit 257fa16f3d6cf6ee17ee984d58eef143db7ecc4c-dirty. Do not edit by hand.
% generated by scripts/paper/tab_detectors.py from results/coverage/coverage.json (57 of 59 clearing cells carry every defense), results/*/detectors/*_metrics.json, results/*/psbd_metrics.json at 2026-09-24T06:01:30.811386+00:00, commit 872fbbbd1f025c8c9924211876e2cda72fc45aab-dirty. Do not edit by hand.
\begin{table}[htbp]
\centering
\small
\caption{TPR at a 10\% false-positive budget of every defense on CIFAR-10. Rows are attacks at their poison rate in percent and columns defenses, PSBD-TM first. Best in each row is bold and second best underlined, gray marks a reading below chance and shaded rows are means. PSBD-RD is our adaptation of the original PSBD site to a ViT block. Scale-Up$^\dagger$ is the data-limited variant and IBD-PSC$^\ast$ the calibrated one.}
\label{tab:detectors-tpr10}
\begin{adjustbox}{max width=\linewidth}
\begin{tabular}{lrrrrrrrrrrrrr}
\toprule
attack, rate \% & \rotatebox{90}{PSBD-TM} & \rotatebox{90}{PSBD-RD} & \rotatebox{90}{Conf} & \rotatebox{90}{STRIP} & \rotatebox{90}{Scale-Up} & \rotatebox{90}{Scale-Up$^\dagger$} & \rotatebox{90}{IBD-PSC} & \rotatebox{90}{IBD-PSC$^\ast$} & \rotatebox{90}{TeCo} & \rotatebox{90}{CD-L} & \rotatebox{90}{Beatrix} & \rotatebox{90}{TED} & \rotatebox{90}{SentiNet} \\
\midrule
BadNets 1 & 0.959 & \textcolor{black!45}{0.015} & \textcolor{black!45}{0.000} & 0.652 & \textbf{0.997} & \textcolor{black!45}{0.000} & \textcolor{black!45}{0.000} & 0.967 & 0.685 & 0.891 & \underline{0.981} & 0.569 & \textcolor{black!45}{0.201} \\
BadNets 5 & 0.965 & \textcolor{black!45}{0.031} & 0.641 & 0.875 & \textbf{1.000} & \textcolor{black!45}{0.000} & \textcolor{black!45}{0.018} & 0.819 & 0.966 & 0.972 & \underline{0.998} & 0.517 & \textcolor{black!45}{0.119} \\
BadNets 10 & 0.935 & \textcolor{black!45}{0.113} & \textcolor{black!45}{0.486} & \textbf{1.000} & \textcolor{black!45}{0.000} & \textcolor{black!45}{0.000} & 0.947 & \underline{0.992} & 0.882 & \textbf{1.000} & \textbf{1.000} & \textcolor{black!45}{0.291} & \textcolor{black!45}{0.222} \\
Blend 1 & 0.773 & 0.983 & \textcolor{black!45}{0.002} & 0.951 & 0.637 & \textcolor{black!45}{0.000} & \textcolor{black!45}{0.140} & \underline{0.998} & 0.764 & 0.996 & \textbf{1.000} & 0.556 & \textcolor{black!45}{0.000} \\
Blend 5 & 0.908 & 0.973 & 0.633 & 0.885 & 0.553 & \textcolor{black!45}{0.000} & 0.670 & 0.997 & 0.689 & \underline{0.998} & \textbf{1.000} & 0.947 & \textcolor{black!45}{0.000} \\
Blend 10 & 0.754 & \underline{0.980} & \textcolor{black!45}{0.004} & 0.853 & \textcolor{black!45}{0.375} & \textcolor{black!45}{0.000} & \textcolor{black!45}{0.000} & \textbf{1.000} & 0.759 & 0.646 & \textbf{1.000} & 0.917 & \textcolor{black!45}{0.008} \\
BPP 1 & 0.956 & 0.897 & \textcolor{black!45}{0.000} & 0.810 & \textcolor{black!45}{0.010} & \textcolor{black!45}{0.010} & \textcolor{black!45}{0.000} & \underline{0.977} & \textcolor{black!45}{0.011} & 0.877 & \textbf{0.985} & 0.642 & \textcolor{black!45}{0.001} \\
BPP 5 & 0.587 & 0.673 & \textcolor{black!45}{0.000} & 0.760 & \textcolor{black!45}{0.010} & \textcolor{black!45}{0.008} & \textcolor{black!45}{0.000} & \underline{0.955} & \textcolor{black!45}{0.014} & 0.856 & \textbf{0.996} & 0.615 & \textcolor{black!45}{0.000} \\
BPP 10 & 0.736 & 0.720 & \textcolor{black!45}{0.013} & 0.881 & \textcolor{black!45}{0.001} & \textcolor{black!45}{0.001} & \textcolor{black!45}{0.041} & \underline{0.990} & \textcolor{black!45}{0.003} & 0.944 & \textbf{0.999} & 0.761 & \textcolor{black!45}{0.001} \\
LF 1 & \textbf{0.954} & 0.753 & \textcolor{black!45}{0.000} & \textcolor{black!45}{0.368} & \textcolor{black!45}{0.379} & \textcolor{black!45}{0.379} & \textcolor{black!45}{0.000} & 0.673 & \textcolor{black!45}{0.133} & \textcolor{black!45}{0.457} & \underline{0.851} & 0.550 & \textcolor{black!45}{0.000} \\
LF 5 & \textbf{0.983} & 0.837 & \textcolor{black!45}{0.000} & 0.867 & 0.762 & \textcolor{black!45}{0.000} & \textcolor{black!45}{0.000} & 0.875 & \textcolor{black!45}{0.033} & 0.722 & \underline{0.970} & 0.502 & \textcolor{black!45}{0.000} \\
LF 10 & \textbf{0.995} & 0.908 & \textcolor{black!45}{0.125} & 0.762 & 0.642 & \textcolor{black!45}{0.000} & \textcolor{black!45}{0.003} & 0.572 & \textcolor{black!45}{0.086} & 0.596 & \underline{0.941} & 0.643 & \textcolor{black!45}{0.001} \\
SIG 10 & \textcolor{black!45}{0.071} & 0.810 & 0.753 & 0.562 & \textcolor{black!45}{0.437} & \textcolor{black!45}{0.437} & \textcolor{black!45}{0.295} & \textcolor{black!45}{0.138} & \textcolor{black!45}{0.296} & \textcolor{black!45}{0.038} & \underline{0.985} & \textbf{0.995} & \textcolor{black!45}{0.002} \\
TaCT 1 & \textcolor{black!45}{0.005} & \textcolor{black!45}{0.013} & \textcolor{black!45}{0.000} & \textcolor{black!45}{0.089} & \textcolor{black!45}{0.131} & \textcolor{black!45}{0.004} & \textcolor{black!45}{0.000} & \textcolor{black!45}{0.134} & \textcolor{black!45}{0.098} & \textcolor{black!45}{0.005} & \textbf{0.994} & \underline{0.853} & \textcolor{black!45}{0.006} \\
TaCT 5 & \textcolor{black!45}{0.000} & \textcolor{black!45}{0.006} & \textcolor{black!45}{0.098} & \textcolor{black!45}{0.141} & 0.823 & 0.691 & \textcolor{black!45}{0.005} & \textcolor{black!45}{0.270} & \textcolor{black!45}{0.310} & \textcolor{black!45}{0.057} & \textbf{0.987} & \underline{0.950} & \textcolor{black!45}{0.000} \\
WaNet 5 & 0.838 & \textbf{0.908} & \textcolor{black!45}{0.000} & \textcolor{black!45}{0.089} & \textcolor{black!45}{0.000} & \textcolor{black!45}{0.005} & \textcolor{black!45}{0.000} & \textcolor{black!45}{0.174} & \underline{0.907} & \textcolor{black!45}{0.378} & \textcolor{black!45}{0.334} & 0.559 & \textcolor{black!45}{0.064} \\
WaNet 10 & \textcolor{black!45}{0.087} & \textbf{0.846} & \textcolor{black!45}{0.000} & \textcolor{black!45}{0.115} & \textcolor{black!45}{0.164} & \textcolor{black!45}{0.060} & \textcolor{black!45}{0.001} & \textcolor{black!45}{0.199} & \textcolor{black!45}{0.361} & \textcolor{black!45}{0.304} & 0.533 & \underline{0.656} & \textcolor{black!45}{0.014} \\
\rowcolor{black!8}\emph{CIFAR-10 mean} & 0.677 & 0.616 & \textcolor{black!45}{0.162} & 0.627 & \textcolor{black!45}{0.407} & \textcolor{black!45}{0.094} & \textcolor{black!45}{0.125} & \underline{0.690} & \textcolor{black!45}{0.412} & 0.631 & \textbf{0.915} & 0.678 & \textcolor{black!45}{0.038} \\
\bottomrule
\end{tabular}
\end{adjustbox}
\end{table}

% generated by scripts/paper/tab_detectors.py from results/coverage/coverage.json (65 of 65 clearing cells carry every defense), results/*/detectors/*_metrics.json, results/*/psbd_metrics.json at 2026-09-24T04:51:48.787049+00:00, commit 9300709d1f6e7a0ba702f153d77765353ffa1368-dirty. Do not edit by hand.
\begin{table}[htbp]
\centering
\small
\caption{TPR at a 10\% false-positive budget of every defense on CIFAR-100. Layout as in \cref{tab:detectors-tpr10}.}
\label{tab:detectors-tpr10-cifar-100}
\begin{adjustbox}{max width=\linewidth}
\begin{tabular}{lrrrrrrrrrrrrr}
\toprule
attack, rate \% & \rotatebox{90}{PSBD-TM} & \rotatebox{90}{PSBD-RD} & \rotatebox{90}{Conf} & \rotatebox{90}{STRIP} & \rotatebox{90}{Scale-Up} & \rotatebox{90}{Scale-Up$^\dagger$} & \rotatebox{90}{IBD-PSC} & \rotatebox{90}{IBD-PSC$^\ast$} & \rotatebox{90}{TeCo} & \rotatebox{90}{CD-L} & \rotatebox{90}{Beatrix} & \rotatebox{90}{TED} & \rotatebox{90}{SentiNet} \\
\midrule
BadNets 1 & 0.999 & \textcolor{black!45}{0.244} & \textcolor{black!45}{0.017} & \textbf{1.000} & 1.000 & \underline{1.000} & \textcolor{black!45}{0.000} & 0.973 & 0.998 & \textbf{1.000} & \textbf{1.000} & 0.975 & 0.932 \\
BadNets 5 & \underline{0.999} & \textcolor{black!45}{0.335} & 0.958 & \textbf{1.000} & \textbf{1.000} & \textbf{1.000} & 0.841 & 0.984 & 0.998 & \textbf{1.000} & \textcolor{black!45}{0.004} & 0.989 & 0.730 \\
BadNets 10 & \textbf{1.000} & \textcolor{black!45}{0.114} & 0.924 & \underline{0.999} & \textbf{1.000} & \textbf{1.000} & 0.994 & 0.987 & 0.999 & \textbf{1.000} & \textbf{1.000} & 0.941 & 0.895 \\
Blend 1 & 0.959 & \textcolor{black!45}{0.347} & \textcolor{black!45}{0.000} & \textcolor{black!45}{0.274} & \textcolor{black!45}{0.335} & \textcolor{black!45}{0.334} & \textcolor{black!45}{0.126} & 0.979 & \underline{0.986} & \textcolor{black!45}{0.423} & \textbf{0.989} & 0.973 & \textcolor{black!45}{0.003} \\
Blend 5 & 0.941 & 0.999 & 0.960 & 0.736 & \textcolor{black!45}{0.266} & \textcolor{black!45}{0.266} & 0.810 & \underline{1.000} & 0.919 & 0.887 & 0.996 & \textbf{1.000} & \textcolor{black!45}{0.459} \\
Blend 10 & \textbf{1.000} & \underline{1.000} & 0.928 & 0.989 & 0.602 & \textcolor{black!45}{0.348} & 0.999 & \textbf{1.000} & \textcolor{black!45}{0.291} & 0.911 & 0.854 & 0.910 & \textcolor{black!45}{0.008} \\
BPP 1 & 0.849 & \textbf{0.914} & \textcolor{black!45}{0.001} & 0.545 & \textcolor{black!45}{0.264} & \textcolor{black!45}{0.008} & \textcolor{black!45}{0.001} & 0.607 & \textcolor{black!45}{0.076} & \underline{0.901} & \textcolor{black!45}{0.038} & 0.898 & \textcolor{black!45}{0.242} \\
BPP 5 & 0.961 & \underline{0.973} & 0.811 & 0.893 & \textcolor{black!45}{0.002} & \textcolor{black!45}{0.001} & 0.523 & 0.824 & 0.752 & \textbf{0.983} & 0.884 & 0.963 & \textcolor{black!45}{0.029} \\
BPP 10 & \underline{0.963} & 0.960 & \textcolor{black!45}{0.486} & 0.860 & \textcolor{black!45}{0.001} & \textcolor{black!45}{0.001} & \textcolor{black!45}{0.260} & 0.948 & \textcolor{black!45}{0.266} & \textbf{0.977} & 0.923 & 0.960 & \textcolor{black!45}{0.022} \\
LF 1 & \textbf{0.942} & 0.903 & \textcolor{black!45}{0.029} & \textcolor{black!45}{0.477} & \textcolor{black!45}{0.233} & \textcolor{black!45}{0.230} & 0.539 & 0.815 & \textcolor{black!45}{0.247} & \textcolor{black!45}{0.049} & 0.887 & \underline{0.917} & \textcolor{black!45}{0.033} \\
LF 5 & \textbf{0.954} & 0.934 & 0.590 & \textcolor{black!45}{0.299} & \textcolor{black!45}{0.117} & \textcolor{black!45}{0.219} & \textcolor{black!45}{0.031} & \textcolor{black!45}{0.243} & \textcolor{black!45}{0.083} & \textcolor{black!45}{0.030} & \textcolor{black!45}{0.366} & \underline{0.949} & \textcolor{black!45}{0.203} \\
LF 10 & \textbf{0.995} & \underline{0.959} & 0.858 & 0.706 & \textcolor{black!45}{0.163} & \textcolor{black!45}{0.162} & 0.765 & 0.920 & \textcolor{black!45}{0.014} & \textcolor{black!45}{0.010} & 0.688 & 0.950 & \textcolor{black!45}{0.001} \\
TaCT 1 & 0.506 & \textcolor{black!45}{0.023} & \underline{0.632} & \textcolor{black!45}{0.184} & \textcolor{black!45}{0.069} & \textcolor{black!45}{0.184} & 0.575 & \textcolor{black!45}{0.218} & \textcolor{black!45}{0.138} & \textcolor{black!45}{0.000} & \textcolor{black!45}{0.000} & \textbf{0.966} & \textcolor{black!45}{0.046} \\
TaCT 5 & 0.540 & \textcolor{black!45}{0.218} & 0.690 & 0.920 & \textbf{0.977} & \textbf{0.977} & 0.586 & 0.736 & 0.552 & \textcolor{black!45}{0.241} & \textcolor{black!45}{0.000} & \underline{0.943} & \textcolor{black!45}{0.080} \\
TaCT 10 & \textcolor{black!45}{0.161} & \textcolor{black!45}{0.000} & 0.724 & \textcolor{black!45}{0.195} & \underline{0.931} & \underline{0.931} & \textcolor{black!45}{0.345} & 0.690 & \textcolor{black!45}{0.379} & \textcolor{black!45}{0.115} & \textcolor{black!45}{0.000} & \textbf{0.943} & \textcolor{black!45}{0.057} \\
\rowcolor{black!8}\emph{CIFAR-100 mean} & \underline{0.851} & 0.595 & 0.574 & 0.672 & \textcolor{black!45}{0.464} & \textcolor{black!45}{0.444} & \textcolor{black!45}{0.493} & 0.795 & 0.513 & 0.569 & 0.575 & \textbf{0.952} & \textcolor{black!45}{0.249} \\
\bottomrule
\end{tabular}
\end{adjustbox}
\end{table}

% generated by scripts/paper/tab_detectors.py from results/coverage/coverage.json (57 of 59 clearing cells carry every defense), results/*/detectors/*_metrics.json, results/*/psbd_metrics.json at 2026-09-24T06:01:44.703285+00:00, commit 872fbbbd1f025c8c9924211876e2cda72fc45aab-dirty. Do not edit by hand.
\begin{table}[htbp]
\centering
\small
\caption{TPR at a 10\% false-positive budget of every defense on GTSRB. Layout as in \cref{tab:detectors-tpr10}.}
\label{tab:detectors-tpr10-gtsrb}
\begin{adjustbox}{max width=\linewidth}
\begin{tabular}{lrrrrrrrrrrrrr}
\toprule
attack, rate \% & \rotatebox{90}{PSBD-TM} & \rotatebox{90}{PSBD-RD} & \rotatebox{90}{Conf} & \rotatebox{90}{STRIP} & \rotatebox{90}{Scale-Up} & \rotatebox{90}{Scale-Up$^\dagger$} & \rotatebox{90}{IBD-PSC} & \rotatebox{90}{IBD-PSC$^\ast$} & \rotatebox{90}{TeCo} & \rotatebox{90}{CD-L} & \rotatebox{90}{Beatrix} & \rotatebox{90}{TED} & \rotatebox{90}{SentiNet} \\
\midrule
BadNets 1 & 0.996 & \textcolor{black!45}{0.001} & \textcolor{black!45}{0.000} & \underline{1.000} & \textcolor{black!45}{0.000} & \textbf{1.000} & \textcolor{black!45}{0.000} & 0.702 & 0.952 & 0.940 & \textbf{1.000} & \textbf{1.000} & \textcolor{black!45}{0.014} \\
BadNets 5 & \textbf{1.000} & 0.659 & \textcolor{black!45}{0.000} & \textbf{1.000} & \textcolor{black!45}{0.000} & \textbf{1.000} & \textcolor{black!45}{0.000} & \textbf{1.000} & \underline{0.999} & \textbf{1.000} & \textbf{1.000} & \textbf{1.000} & \textcolor{black!45}{0.035} \\
BadNets 10 & \textbf{1.000} & 0.999 & \textcolor{black!45}{0.000} & \textbf{1.000} & \textcolor{black!45}{0.000} & \textbf{1.000} & \textcolor{black!45}{0.000} & \textbf{1.000} & 0.862 & \underline{1.000} & \textbf{1.000} & \textbf{1.000} & \textcolor{black!45}{0.001} \\
Blend 1 & \textbf{1.000} & \underline{1.000} & \textcolor{black!45}{0.000} & 0.935 & \textcolor{black!45}{0.000} & 0.542 & \textcolor{black!45}{0.000} & \textbf{1.000} & \textcolor{black!45}{0.388} & 0.742 & \textbf{1.000} & \textbf{1.000} & \textcolor{black!45}{0.000} \\
Blend 5 & \underline{1.000} & 0.714 & 0.727 & 0.953 & \textcolor{black!45}{0.000} & 0.867 & 0.912 & 0.994 & 0.750 & 0.961 & \textbf{1.000} & \underline{1.000} & \textcolor{black!45}{0.048} \\
Blend 10 & \underline{1.000} & 0.996 & \textcolor{black!45}{0.000} & 0.936 & \textcolor{black!45}{0.000} & 0.874 & 1.000 & \underline{1.000} & \textcolor{black!45}{0.436} & 1.000 & \textbf{1.000} & \underline{1.000} & \textcolor{black!45}{0.141} \\
BPP 1 & 0.717 & \textcolor{black!45}{0.128} & \textcolor{black!45}{0.000} & \textcolor{black!45}{0.006} & \textcolor{black!45}{0.000} & \textcolor{black!45}{0.005} & \textcolor{black!45}{0.000} & 0.873 & \textcolor{black!45}{0.154} & 0.826 & \textbf{0.964} & \underline{0.895} & \textcolor{black!45}{0.056} \\
BPP 5 & 0.978 & 0.988 & \textcolor{black!45}{0.000} & 0.802 & \textcolor{black!45}{0.000} & \textcolor{black!45}{0.001} & \textcolor{black!45}{0.000} & \textcolor{black!45}{0.438} & \textcolor{black!45}{0.026} & 0.965 & \textbf{0.994} & \underline{0.992} & \textcolor{black!45}{0.067} \\
BPP 10 & 0.971 & 0.971 & \textcolor{black!45}{0.000} & 0.762 & \textcolor{black!45}{0.000} & \textcolor{black!45}{0.000} & \textcolor{black!45}{0.000} & \textcolor{black!45}{0.218} & \textcolor{black!45}{0.066} & 0.919 & \underline{0.995} & \textbf{0.996} & \textcolor{black!45}{0.007} \\
LF 1 & 0.926 & 0.977 & \textcolor{black!45}{0.000} & 0.694 & \textcolor{black!45}{0.000} & 0.830 & \textcolor{black!45}{0.000} & 0.906 & \textcolor{black!45}{0.237} & \textcolor{black!45}{0.047} & \textbf{0.983} & \underline{0.979} & \textcolor{black!45}{0.002} \\
LF 5 & 0.999 & 0.997 & \textcolor{black!45}{0.000} & 0.915 & \textcolor{black!45}{0.000} & 0.781 & \textcolor{black!45}{0.000} & 0.937 & \textcolor{black!45}{0.026} & \textcolor{black!45}{0.075} & \textbf{0.999} & \underline{0.999} & \textcolor{black!45}{0.000} \\
LF 10 & 0.998 & 0.999 & \textcolor{black!45}{0.000} & 0.928 & \textcolor{black!45}{0.000} & 0.750 & \textcolor{black!45}{0.000} & 0.873 & \textcolor{black!45}{0.034} & \textcolor{black!45}{0.076} & \textbf{1.000} & \underline{0.999} & \textcolor{black!45}{0.000} \\
TaCT 5 & 0.502 & \textcolor{black!45}{0.010} & \textcolor{black!45}{0.000} & \textcolor{black!45}{0.426} & \textcolor{black!45}{0.000} & \textbf{1.000} & \textcolor{black!45}{0.094} & 0.724 & \underline{0.993} & \textcolor{black!45}{0.367} & \textbf{1.000} & \textbf{1.000} & \textcolor{black!45}{0.003} \\
WaNet 10 & 0.949 & \underline{0.950} & \textcolor{black!45}{0.000} & \textcolor{black!45}{0.077} & \textcolor{black!45}{0.000} & \textcolor{black!45}{0.210} & \textcolor{black!45}{0.060} & \textcolor{black!45}{0.220} & 0.925 & \textcolor{black!45}{0.056} & \textcolor{black!45}{0.044} & \textbf{0.953} & \textcolor{black!45}{0.005} \\
\rowcolor{black!8}\emph{GTSRB mean} & \underline{0.931} & 0.742 & \textcolor{black!45}{0.052} & 0.745 & \textcolor{black!45}{0.000} & 0.633 & \textcolor{black!45}{0.148} & 0.778 & \textcolor{black!45}{0.489} & 0.641 & 0.927 & \textbf{0.987} & \textcolor{black!45}{0.027} \\
\bottomrule
\end{tabular}
\end{adjustbox}
\end{table}

% generated by scripts/paper/tab_detectors.py from results/coverage/coverage.json (57 of 59 clearing cells carry every defense), results/*/detectors/*_metrics.json, results/*/psbd_metrics.json at 2026-09-29T16:27:10.151649+00:00, commit 19488b81358b06040ba96f361d5061b2981f1f25-dirty. Do not edit by hand.
\begin{table}[htbp]
\centering
\small
\caption{TPR at a 10\% false-positive budget of every defense on Tiny ImageNet. Layout as in \cref{tab:detectors-tpr10}.}
\label{tab:detectors-tpr10-tiny-imagenet}
\begin{adjustbox}{max width=\linewidth}
\begin{tabular}{lrrrrrrrrrrrrr}
\toprule
attack, rate \% & \rotatebox{90}{PSBD-TM} & \rotatebox{90}{PSBD-RD} & \rotatebox{90}{Conf} & \rotatebox{90}{STRIP} & \rotatebox{90}{Scale-Up} & \rotatebox{90}{Scale-Up$^\dagger$} & \rotatebox{90}{IBD-PSC} & \rotatebox{90}{IBD-PSC$^\ast$} & \rotatebox{90}{TeCo} & \rotatebox{90}{CD-L} & \rotatebox{90}{Beatrix} & \rotatebox{90}{TED} & \rotatebox{90}{SentiNet} \\
\midrule
BadNets 1 & 0.989 & 0.700 & \textcolor{black!45}{0.003} & 0.769 & \textbf{0.997} & \textbf{0.997} & 0.758 & 0.916 & \underline{0.992} & \textcolor{black!45}{0.000} & \textcolor{black!45}{0.002} & 0.802 & \textcolor{black!45}{0.192} \\
BadNets 5 & \underline{0.999} & 0.566 & 0.800 & 0.868 & \textbf{0.999} & \underline{0.999} & 0.828 & 0.871 & 0.991 & \textcolor{black!45}{0.000} & \textcolor{black!45}{0.000} & 0.549 & \textcolor{black!45}{0.062} \\
BadNets 10 & \underline{1.000} & 0.951 & 0.960 & 0.769 & \textbf{1.000} & 1.000 & 0.877 & 0.929 & 0.987 & \textcolor{black!45}{0.000} & \textcolor{black!45}{0.000} & \textcolor{black!45}{0.487} & \textcolor{black!45}{0.003} \\
Blend 1 & 0.985 & \underline{0.994} & 0.541 & 0.634 & 0.638 & \textcolor{black!45}{0.273} & 0.966 & \textbf{0.999} & 0.709 & \textcolor{black!45}{0.153} & \textcolor{black!45}{0.001} & \textcolor{black!45}{0.064} & \textcolor{black!45}{0.001} \\
Blend 5 & 0.997 & \underline{0.999} & 0.931 & 0.808 & 0.741 & \textcolor{black!45}{0.388} & 0.998 & \textbf{0.999} & 0.701 & \textcolor{black!45}{0.063} & \textcolor{black!45}{0.000} & \textcolor{black!45}{0.007} & \textcolor{black!45}{0.001} \\
Blend 10 & 0.999 & \underline{0.999} & 0.990 & 0.940 & 0.879 & \textcolor{black!45}{0.000} & 0.999 & \textbf{1.000} & 0.852 & \textcolor{black!45}{0.000} & \textcolor{black!45}{0.000} & \textcolor{black!45}{0.060} & \textcolor{black!45}{0.001} \\
BPP 1 & \textbf{0.972} & \underline{0.972} & \textcolor{black!45}{0.382} & 0.907 & \textcolor{black!45}{0.031} & \textcolor{black!45}{0.029} & 0.959 & 0.970 & \textcolor{black!45}{0.217} & \textcolor{black!45}{0.009} & \textcolor{black!45}{0.002} & \textcolor{black!45}{0.231} & \textcolor{black!45}{0.221} \\
BPP 5 & \textbf{0.984} & 0.977 & 0.779 & 0.873 & \textcolor{black!45}{0.003} & \textcolor{black!45}{0.003} & 0.957 & \underline{0.981} & \textcolor{black!45}{0.100} & \textcolor{black!45}{0.122} & \textcolor{black!45}{0.157} & 0.597 & \textcolor{black!45}{0.312} \\
BPP 10 & \underline{0.981} & \textbf{0.995} & \textcolor{black!45}{0.451} & 0.925 & \textcolor{black!45}{0.007} & \textcolor{black!45}{0.001} & 0.566 & 0.566 & \textcolor{black!45}{0.150} & \textcolor{black!45}{0.053} & \textcolor{black!45}{0.007} & \textcolor{black!45}{0.055} & \textcolor{black!45}{0.122} \\
LF 1 & \textbf{0.929} & \underline{0.918} & \textcolor{black!45}{0.448} & \textcolor{black!45}{0.222} & \textcolor{black!45}{0.430} & \textcolor{black!45}{0.324} & 0.734 & 0.734 & \textcolor{black!45}{0.285} & \textcolor{black!45}{0.283} & \textcolor{black!45}{0.004} & \textcolor{black!45}{0.263} & \textcolor{black!45}{0.011} \\
LF 5 & \textbf{0.987} & \underline{0.987} & 0.893 & 0.617 & 0.565 & \textcolor{black!45}{0.001} & 0.954 & 0.954 & \textcolor{black!45}{0.141} & \textcolor{black!45}{0.011} & \textcolor{black!45}{0.000} & \textcolor{black!45}{0.039} & \textcolor{black!45}{0.003} \\
LF 10 & \textbf{0.973} & \underline{0.970} & 0.838 & \textcolor{black!45}{0.454} & \textcolor{black!45}{0.404} & \textcolor{black!45}{0.404} & 0.910 & 0.910 & \textcolor{black!45}{0.361} & \textcolor{black!45}{0.034} & \textcolor{black!45}{0.002} & \textcolor{black!45}{0.152} & \textcolor{black!45}{0.003} \\
WaNet 5 & \textbf{0.924} & \underline{0.922} & \textcolor{black!45}{0.401} & \textcolor{black!45}{0.105} & \textcolor{black!45}{0.055} & \textcolor{black!45}{0.058} & 0.680 & 0.680 & 0.841 & \textcolor{black!45}{0.070} & \textcolor{black!45}{0.003} & \textcolor{black!45}{0.052} & \textcolor{black!45}{0.011} \\
WaNet 10 & \textbf{0.932} & \underline{0.924} & 0.611 & \textcolor{black!45}{0.112} & \textcolor{black!45}{0.130} & \textcolor{black!45}{0.059} & 0.677 & 0.677 & 0.890 & 0.742 & \textcolor{black!45}{0.001} & \textcolor{black!45}{0.292} & \textcolor{black!45}{0.011} \\
\rowcolor{black!8}\emph{Tiny ImageNet mean} & \textbf{0.975} & \underline{0.920} & 0.645 & 0.643 & \textcolor{black!45}{0.491} & \textcolor{black!45}{0.324} & 0.847 & 0.870 & 0.587 & \textcolor{black!45}{0.110} & \textcolor{black!45}{0.013} & \textcolor{black!45}{0.261} & \textcolor{black!45}{0.068} \\
\rowcolor{black!8}\emph{all models} & \textbf{0.873} & 0.744 & \textcolor{black!45}{0.335} & 0.682 & \textcolor{black!45}{0.329} & \textcolor{black!45}{0.343} & \textcolor{black!45}{0.385} & \underline{0.791} & 0.503 & 0.516 & 0.655 & 0.709 & \textcolor{black!45}{0.097} \\
\bottomrule
\end{tabular}
\end{adjustbox}
\end{table}


\FloatBarrier
% generated by scripts/paper/tab_detectors.py from results/coverage/coverage.json (54 of 56 successful cells carry every defense), results/*/detectors/*_metrics.json, results/*/psbd_metrics.json at 2026-09-29T18:07:02.827677+00:00, commit b2d32cf11708d5de965d3e13604c863a3ad9b493-dirty. Do not edit by hand.
% generated by scripts/paper/tab_detectors.py from results/coverage/coverage.json (65 of 65 clearing cells carry every defense), results/*/detectors/*_metrics.json, results/*/psbd_metrics.json at 2026-09-24T04:52:19.118810+00:00, commit 9300709d1f6e7a0ba702f153d77765353ffa1368-dirty. Do not edit by hand.
\begin{table}[htbp]
\centering
\small
\caption{TPR at a 20\% false-positive budget of every defense on CIFAR-10. Rows are attacks at their poison rate in percent and columns defenses, PSBD-TM first. Best in each row is bold and second best underlined, gray marks a reading below chance and shaded rows are means. PSBD-RD is our adaptation of the original PSBD site to a ViT block. Scale-Up$^\dagger$ is the data-limited variant and IBD-PSC$^\ast$ the calibrated one.}
\label{tab:detectors-tpr20}
\begin{adjustbox}{max width=\linewidth}
\begin{tabular}{lrrrrrrrrrrrrr}
\toprule
attack, rate \% & \rotatebox{90}{PSBD-TM} & \rotatebox{90}{PSBD-RD} & \rotatebox{90}{Conf} & \rotatebox{90}{STRIP} & \rotatebox{90}{Scale-Up} & \rotatebox{90}{Scale-Up$^\dagger$} & \rotatebox{90}{IBD-PSC} & \rotatebox{90}{IBD-PSC$^\ast$} & \rotatebox{90}{TeCo} & \rotatebox{90}{CD-L} & \rotatebox{90}{Beatrix} & \rotatebox{90}{TED} & \rotatebox{90}{SentiNet} \\
\midrule
BadNets 1 & 0.978 & \textcolor{black!45}{0.046} & \textcolor{black!45}{0.000} & 0.842 & \underline{0.997} & \underline{0.997} & \textcolor{black!45}{0.000} & 0.985 & 0.879 & 0.914 & \textbf{1.000} & 0.698 & \textcolor{black!45}{0.437} \\
BadNets 5 & 0.995 & \textcolor{black!45}{0.168} & 0.784 & 0.989 & \underline{1.000} & \underline{1.000} & \textcolor{black!45}{0.138} & 0.964 & 0.978 & 0.983 & \textbf{1.000} & 0.672 & \textcolor{black!45}{0.178} \\
BadNets 10 & \underline{1.000} & \textcolor{black!45}{0.301} & 0.792 & \textbf{1.000} & \textbf{1.000} & \textcolor{black!45}{0.000} & 0.998 & \textbf{1.000} & 0.933 & \textbf{1.000} & \textbf{1.000} & \textcolor{black!45}{0.447} & \textcolor{black!45}{0.278} \\
Blend 1 & 0.869 & 0.992 & \textcolor{black!45}{0.109} & 0.979 & 0.767 & 0.523 & 0.934 & \underline{1.000} & 0.835 & 0.998 & \textbf{1.000} & 0.716 & \textcolor{black!45}{0.000} \\
Blend 5 & 0.989 & 0.995 & 0.891 & 0.948 & 0.777 & 0.553 & 0.895 & \underline{1.000} & 0.831 & 0.999 & \textbf{1.000} & 0.979 & \textcolor{black!45}{0.000} \\
Blend 10 & 0.849 & \underline{0.992} & \textcolor{black!45}{0.163} & 0.940 & 0.525 & \textcolor{black!45}{0.375} & \textcolor{black!45}{0.095} & \textbf{1.000} & 0.817 & 0.930 & \textbf{1.000} & 0.978 & \textcolor{black!45}{0.035} \\
BPP 1 & 0.976 & 0.952 & \textcolor{black!45}{0.000} & 0.886 & \textcolor{black!45}{0.035} & \textcolor{black!45}{0.025} & \textcolor{black!45}{0.000} & \underline{0.980} & \textcolor{black!45}{0.060} & 0.916 & \textbf{0.989} & 0.787 & \textcolor{black!45}{0.008} \\
BPP 5 & 0.655 & 0.788 & \textcolor{black!45}{0.002} & 0.854 & \textcolor{black!45}{0.023} & \textcolor{black!45}{0.019} & \textcolor{black!45}{0.001} & \underline{0.978} & \textcolor{black!45}{0.071} & 0.895 & \textbf{0.998} & 0.786 & \textcolor{black!45}{0.002} \\
BPP 10 & 0.830 & 0.799 & \textcolor{black!45}{0.379} & 0.931 & \textcolor{black!45}{0.002} & \textcolor{black!45}{0.002} & \textcolor{black!45}{0.499} & \underline{0.991} & \textcolor{black!45}{0.069} & 0.955 & \textbf{0.999} & 0.838 & \textcolor{black!45}{0.004} \\
LF 1 & \textbf{0.975} & 0.898 & \textcolor{black!45}{0.000} & 0.555 & 0.556 & \textcolor{black!45}{0.462} & \textcolor{black!45}{0.000} & 0.814 & \textcolor{black!45}{0.213} & 0.637 & \underline{0.938} & 0.669 & \textcolor{black!45}{0.005} \\
LF 5 & \textbf{0.994} & 0.947 & \textcolor{black!45}{0.000} & 0.919 & 0.847 & 0.666 & \textcolor{black!45}{0.000} & 0.962 & \textcolor{black!45}{0.057} & 0.855 & \underline{0.992} & 0.638 & \textcolor{black!45}{0.002} \\
LF 10 & \textbf{0.999} & 0.972 & 0.553 & 0.835 & 0.839 & 0.567 & \textcolor{black!45}{0.062} & 0.782 & \textcolor{black!45}{0.125} & 0.838 & \underline{0.983} & 0.756 & \textcolor{black!45}{0.006} \\
SIG 10 & \textcolor{black!45}{0.097} & 0.878 & 0.804 & 0.691 & 0.541 & 0.541 & \textcolor{black!45}{0.327} & \textcolor{black!45}{0.307} & 0.615 & \textcolor{black!45}{0.243} & \underline{0.993} & \textbf{0.997} & \textcolor{black!45}{0.003} \\
TaCT 1 & \textcolor{black!45}{0.009} & \textcolor{black!45}{0.038} & \textcolor{black!45}{0.000} & \textcolor{black!45}{0.207} & \textcolor{black!45}{0.279} & \textcolor{black!45}{0.133} & \textcolor{black!45}{0.000} & \textcolor{black!45}{0.441} & \textcolor{black!45}{0.246} & \textcolor{black!45}{0.064} & \textbf{1.000} & \underline{0.908} & \textcolor{black!45}{0.009} \\
TaCT 5 & \textcolor{black!45}{0.031} & \textcolor{black!45}{0.020} & \textcolor{black!45}{0.317} & \textcolor{black!45}{0.270} & 0.907 & 0.823 & \textcolor{black!45}{0.074} & 0.533 & 0.569 & \textcolor{black!45}{0.285} & \textbf{1.000} & \underline{0.971} & \textcolor{black!45}{0.003} \\
TaCT 10 & \textcolor{black!45}{0.206} & \textcolor{black!45}{0.000} & \textcolor{black!45}{0.152} & 0.575 & \underline{0.992} & \underline{0.992} & \textcolor{black!45}{0.339} & 0.872 & \textcolor{black!45}{0.389} & 0.799 & \textbf{1.000} & 0.908 & \textcolor{black!45}{0.011} \\
WaNet 5 & 0.913 & \textbf{0.965} & \textcolor{black!45}{0.000} & \textcolor{black!45}{0.182} & 0.889 & \textcolor{black!45}{0.008} & \textcolor{black!45}{0.002} & \textcolor{black!45}{0.350} & \underline{0.952} & 0.502 & \textcolor{black!45}{0.491} & 0.689 & \textcolor{black!45}{0.130} \\
WaNet 10 & \textcolor{black!45}{0.268} & \textbf{0.914} & \textcolor{black!45}{0.006} & \textcolor{black!45}{0.215} & \textcolor{black!45}{0.274} & \textcolor{black!45}{0.176} & \textcolor{black!45}{0.027} & \textcolor{black!45}{0.379} & 0.529 & \textcolor{black!45}{0.371} & 0.728 & \underline{0.754} & \textcolor{black!45}{0.036} \\
\rowcolor{black!8}\emph{CIFAR-10 mean} & 0.702 & 0.648 & \textcolor{black!45}{0.275} & 0.712 & 0.625 & \textcolor{black!45}{0.437} & \textcolor{black!45}{0.244} & \underline{0.797} & 0.509 & 0.732 & \textbf{0.951} & 0.788 & \textcolor{black!45}{0.064} \\
\bottomrule
\end{tabular}
\end{adjustbox}
\end{table}

% generated by scripts/paper/tab_detectors.py from results/coverage/coverage.json (65 of 65 clearing cells carry every defense), results/*/detectors/*_metrics.json, results/*/psbd_metrics.json at 2026-09-24T04:39:51.453428+00:00, commit 52e4ba1a7ac941365bc37d68e1e902dad4fb4f74-dirty. Do not edit by hand.
\begin{table}[htbp]
\centering
\small
\caption{TPR at a 20\% false-positive budget of every defense on CIFAR-100. Layout as in \cref{tab:detectors-tpr20}.}
\label{tab:detectors-tpr20-cifar-100}
\begin{adjustbox}{max width=\linewidth}
\begin{tabular}{lrrrrrrrrrrrrr}
\toprule
attack, rate \% & \rotatebox{90}{PSBD-TM} & \rotatebox{90}{PSBD-RD} & \rotatebox{90}{Conf} & \rotatebox{90}{STRIP} & \rotatebox{90}{Scale-Up} & \rotatebox{90}{Scale-Up$^\dagger$} & \rotatebox{90}{IBD-PSC} & \rotatebox{90}{IBD-PSC$^\ast$} & \rotatebox{90}{TeCo} & \rotatebox{90}{CD-L} & \rotatebox{90}{Beatrix} & \rotatebox{90}{TED} & \rotatebox{90}{SentiNet} \\
\midrule
BadNets 1 & \underline{1.000} & \textcolor{black!45}{0.476} & \textcolor{black!45}{0.305} & \textbf{1.000} & \underline{1.000} & \textbf{1.000} & 0.633 & 1.000 & 1.000 & \textbf{1.000} & \textbf{1.000} & 0.985 & 0.959 \\
BadNets 5 & \textbf{1.000} & 0.672 & 0.986 & \textbf{1.000} & \textbf{1.000} & \textbf{1.000} & 0.933 & 0.997 & \underline{0.999} & \textbf{1.000} & \textcolor{black!45}{0.102} & 0.994 & 0.806 \\
BadNets 10 & \textbf{1.000} & \textcolor{black!45}{0.234} & 0.973 & \textbf{1.000} & \textbf{1.000} & \textbf{1.000} & \textbf{1.000} & \underline{0.999} & \textbf{1.000} & \textbf{1.000} & \textbf{1.000} & 0.979 & 0.936 \\
Blend 1 & 0.983 & 0.598 & \textcolor{black!45}{0.003} & \textcolor{black!45}{0.450} & \textcolor{black!45}{0.335} & \textcolor{black!45}{0.334} & 0.734 & 0.988 & \textbf{0.991} & 0.774 & \underline{0.989} & 0.983 & \textcolor{black!45}{0.010} \\
Blend 5 & 0.976 & \underline{1.000} & 0.995 & 0.861 & \textcolor{black!45}{0.470} & \textcolor{black!45}{0.470} & 0.982 & \underline{1.000} & 0.986 & 0.991 & \underline{1.000} & \textbf{1.000} & 0.884 \\
Blend 10 & \textbf{1.000} & \textbf{1.000} & 0.996 & 0.993 & 0.793 & 0.602 & \underline{1.000} & \textbf{1.000} & 0.845 & 0.999 & 0.991 & 0.996 & \textcolor{black!45}{0.070} \\
BPP 1 & 0.919 & \underline{0.929} & \textcolor{black!45}{0.002} & 0.716 & \textcolor{black!45}{0.276} & \textcolor{black!45}{0.121} & \textcolor{black!45}{0.012} & 0.842 & \textcolor{black!45}{0.160} & \textbf{0.942} & 0.635 & 0.913 & 0.787 \\
BPP 5 & 0.977 & \underline{0.984} & 0.893 & 0.934 & \textcolor{black!45}{0.328} & \textcolor{black!45}{0.004} & 0.865 & 0.957 & 0.851 & \textbf{0.989} & 0.954 & 0.977 & 0.707 \\
BPP 10 & 0.975 & 0.979 & 0.853 & 0.920 & \textcolor{black!45}{0.294} & \textcolor{black!45}{0.003} & 0.864 & 0.984 & \textcolor{black!45}{0.488} & \underline{0.985} & \textbf{0.986} & 0.975 & \textcolor{black!45}{0.064} \\
LF 1 & \textbf{0.947} & 0.940 & \textcolor{black!45}{0.189} & 0.616 & \textcolor{black!45}{0.391} & \textcolor{black!45}{0.388} & 0.728 & 0.917 & \textcolor{black!45}{0.374} & \textcolor{black!45}{0.101} & \underline{0.946} & 0.937 & \textcolor{black!45}{0.101} \\
LF 5 & \textbf{0.959} & 0.955 & 0.693 & \textcolor{black!45}{0.418} & \textcolor{black!45}{0.221} & 0.505 & \textcolor{black!45}{0.146} & 0.521 & \textcolor{black!45}{0.191} & \textcolor{black!45}{0.054} & 0.954 & \underline{0.959} & 0.788 \\
LF 10 & \textbf{0.995} & 0.983 & 0.933 & 0.815 & \textcolor{black!45}{0.163} & \textcolor{black!45}{0.480} & 0.879 & 0.978 & \textcolor{black!45}{0.032} & \textcolor{black!45}{0.033} & \underline{0.995} & 0.976 & \textcolor{black!45}{0.263} \\
TaCT 1 & 0.690 & \textcolor{black!45}{0.103} & 0.736 & \textcolor{black!45}{0.299} & \textcolor{black!45}{0.184} & 0.506 & \underline{0.770} & 0.667 & \textcolor{black!45}{0.287} & \textcolor{black!45}{0.126} & \textcolor{black!45}{0.000} & \textbf{0.977} & \textcolor{black!45}{0.149} \\
TaCT 5 & 0.862 & 0.552 & 0.828 & 0.966 & \textbf{0.989} & \textbf{0.989} & 0.782 & 0.908 & 0.805 & 0.805 & \textcolor{black!45}{0.000} & \underline{0.977} & \textcolor{black!45}{0.241} \\
TaCT 10 & 0.586 & \textcolor{black!45}{0.103} & 0.816 & \textcolor{black!45}{0.425} & \underline{0.943} & \underline{0.943} & 0.747 & 0.931 & 0.644 & 0.552 & \textcolor{black!45}{0.000} & \textbf{0.977} & \textcolor{black!45}{0.230} \\
\rowcolor{black!8}\emph{CIFAR-100 mean} & \underline{0.925} & 0.701 & 0.680 & 0.761 & 0.559 & 0.556 & 0.738 & 0.912 & 0.643 & 0.690 & 0.703 & \textbf{0.974} & \textcolor{black!45}{0.466} \\
\bottomrule
\end{tabular}
\end{adjustbox}
\end{table}

% generated by scripts/paper/tab_detectors.py from results/coverage/coverage.json (56 of 62 successful cells carry every defense), results/*/detectors/*_metrics.json, results/*/psbd_metrics.json at 2026-09-30T14:56:51.016962+00:00, commit 1fd38381a58031e98490ed24679d373b446725fb-dirty. Do not edit by hand.
\begin{table}[htbp]
\centering
\small
\caption{TPR at a 20\% false-positive budget of every defense on GTSRB. Layout as in \cref{tab:detectors-tpr20}.}
\label{tab:detectors-tpr20-gtsrb}
\begin{adjustbox}{max width=\linewidth}
\begin{tabular}{lrrrrrrrrrrrrr}
\toprule
attack, rate \% & \rotatebox{90}{PSBD-TM} & \rotatebox{90}{PSBD-RD} & \rotatebox{90}{Conf} & \rotatebox{90}{STRIP} & \rotatebox{90}{Scale-Up} & \rotatebox{90}{Scale-Up$^\dagger$} & \rotatebox{90}{IBD-PSC} & \rotatebox{90}{IBD-PSC$^\ast$} & \rotatebox{90}{TeCo} & \rotatebox{90}{CD-L} & \rotatebox{90}{Beatrix} & \rotatebox{90}{TED} & \rotatebox{90}{SentiNet} \\
\midrule
BadNets 1 & \underline{0.999} & \textcolor{black!45}{0.008} & \textcolor{black!45}{0.000} & \textbf{1.000} & \textbf{1.000} & \textbf{1.000} & \textcolor{black!45}{0.003} & 0.936 & 0.988 & 0.992 & \textbf{1.000} & \textbf{1.000} & \textcolor{black!45}{0.014} \\
BadNets 5 & \textbf{1.000} & \underline{0.830} & \textcolor{black!45}{0.075} & \textbf{1.000} & \textbf{1.000} & \textbf{1.000} & \textcolor{black!45}{0.000} & \textbf{1.000} & \textbf{1.000} & \textbf{1.000} & \textbf{1.000} & \textbf{1.000} & \textcolor{black!45}{0.042} \\
BadNets 10 & \textbf{1.000} & \underline{1.000} & \textcolor{black!45}{0.000} & \textbf{1.000} & \textbf{1.000} & \textbf{1.000} & \textcolor{black!45}{0.000} & \textbf{1.000} & 0.960 & \textbf{1.000} & \textbf{1.000} & \textbf{1.000} & \textcolor{black!45}{0.001} \\
Blend 1 & \underline{1.000} & \underline{1.000} & \textcolor{black!45}{0.000} & 0.949 & \textcolor{black!45}{0.000} & 0.718 & \textcolor{black!45}{0.000} & \underline{1.000} & 0.528 & 0.963 & \textbf{1.000} & \underline{1.000} & \textcolor{black!45}{0.000} \\
Blend 5 & \underline{1.000} & 0.834 & 0.844 & 0.966 & \textcolor{black!45}{0.361} & 0.867 & 0.952 & 0.999 & 0.861 & 0.988 & \textbf{1.000} & \underline{1.000} & \textcolor{black!45}{0.064} \\
Blend 10 & \underline{1.000} & 0.998 & \textcolor{black!45}{0.000} & 0.947 & \textcolor{black!45}{0.374} & 0.874 & 1.000 & \underline{1.000} & 0.643 & \textbf{1.000} & \textbf{1.000} & \underline{1.000} & \textcolor{black!45}{0.157} \\
BPP 1 & 0.809 & 0.507 & \textcolor{black!45}{0.000} & \textcolor{black!45}{0.010} & \textcolor{black!45}{0.029} & \textcolor{black!45}{0.017} & \textcolor{black!45}{0.000} & 0.881 & \textcolor{black!45}{0.247} & \underline{0.906} & \textbf{0.983} & 0.903 & \textcolor{black!45}{0.145} \\
BPP 5 & 0.987 & 0.992 & \textcolor{black!45}{0.043} & 0.855 & \textcolor{black!45}{0.001} & \textcolor{black!45}{0.008} & \textcolor{black!45}{0.000} & 0.624 & \textcolor{black!45}{0.067} & 0.981 & \textbf{0.998} & \underline{0.993} & \textcolor{black!45}{0.078} \\
BPP 10 & 0.986 & 0.992 & \textcolor{black!45}{0.002} & 0.824 & \textcolor{black!45}{0.000} & \textcolor{black!45}{0.000} & \textcolor{black!45}{0.000} & \textcolor{black!45}{0.440} & \textcolor{black!45}{0.116} & 0.980 & \textbf{0.998} & \underline{0.996} & \textcolor{black!45}{0.011} \\
Label-Consistent 5 & 0.963 & \textcolor{black!45}{0.088} & \textcolor{black!45}{0.015} & 0.799 & 0.931 & 0.958 & \textcolor{black!45}{0.051} & 0.507 & 0.825 & \textcolor{black!45}{0.422} & \textbf{0.999} & 0.963 & \underline{0.975} \\
LF 1 & 0.973 & \underline{0.985} & \textcolor{black!45}{0.000} & 0.773 & \textcolor{black!45}{0.000} & 0.830 & \textcolor{black!45}{0.001} & 0.928 & \textcolor{black!45}{0.337} & \textcolor{black!45}{0.075} & \textbf{0.988} & 0.982 & \textcolor{black!45}{0.004} \\
LF 5 & \underline{0.999} & \underline{0.999} & \textcolor{black!45}{0.000} & 0.928 & 0.645 & 0.862 & \textcolor{black!45}{0.000} & 0.971 & \textcolor{black!45}{0.038} & \textcolor{black!45}{0.263} & \textbf{1.000} & 0.999 & \textcolor{black!45}{0.000} \\
LF 10 & 0.999 & \underline{0.999} & \textcolor{black!45}{0.000} & 0.939 & \textcolor{black!45}{0.474} & 0.750 & \textcolor{black!45}{0.072} & 0.921 & \textcolor{black!45}{0.052} & \textcolor{black!45}{0.177} & \textbf{1.000} & \textbf{1.000} & \textcolor{black!45}{0.000} \\
TaCT 1 & \textcolor{black!45}{0.000} & 0.813 & \textcolor{black!45}{0.000} & \textcolor{black!45}{0.291} & \textcolor{black!45}{0.000} & 0.585 & \textcolor{black!45}{0.012} & \underline{0.853} & 0.744 & \textcolor{black!45}{0.104} & \textbf{1.000} & \textbf{1.000} & \textcolor{black!45}{0.071} \\
TaCT 5 & 0.818 & \textcolor{black!45}{0.041} & \textcolor{black!45}{0.221} & 0.603 & \textcolor{black!45}{0.000} & \textbf{1.000} & \textcolor{black!45}{0.468} & \underline{0.904} & \textbf{1.000} & 0.719 & \textbf{1.000} & \textbf{1.000} & \textcolor{black!45}{0.003} \\
\rowcolor{black!8}\emph{GTSRB mean} & 0.902 & 0.739 & \textcolor{black!45}{0.080} & 0.792 & \textcolor{black!45}{0.388} & 0.698 & \textcolor{black!45}{0.171} & 0.864 & 0.560 & 0.705 & \textbf{0.998} & \underline{0.989} & \textcolor{black!45}{0.104} \\
\bottomrule
\end{tabular}
\end{adjustbox}
\end{table}

% generated by scripts/paper/tab_detectors.py from results/coverage/coverage.json (57 of 59 clearing cells carry every defense), results/*/detectors/*_metrics.json, results/*/psbd_metrics.json at 2026-09-24T06:02:32.914232+00:00, commit 872fbbbd1f025c8c9924211876e2cda72fc45aab-dirty. Do not edit by hand.
\begin{table}[htbp]
\centering
\small
\caption{TPR at a 20\% false-positive budget of every defense on Tiny ImageNet. Layout as in \cref{tab:detectors-tpr20}.}
\label{tab:detectors-tpr20-tiny-imagenet}
\begin{adjustbox}{max width=\linewidth}
\begin{tabular}{lrrrrrrrrrrrrr}
\toprule
attack, rate \% & \rotatebox{90}{PSBD-TM} & \rotatebox{90}{PSBD-RD} & \rotatebox{90}{Conf} & \rotatebox{90}{STRIP} & \rotatebox{90}{Scale-Up} & \rotatebox{90}{Scale-Up$^\dagger$} & \rotatebox{90}{IBD-PSC} & \rotatebox{90}{IBD-PSC$^\ast$} & \rotatebox{90}{TeCo} & \rotatebox{90}{CD-L} & \rotatebox{90}{Beatrix} & \rotatebox{90}{TED} & \rotatebox{90}{SentiNet} \\
\midrule
BadNets 1 & \underline{0.997} & 0.836 & \textcolor{black!45}{0.243} & 0.937 & \underline{0.997} & \underline{0.997} & 0.980 & 0.979 & \textbf{0.998} & \textcolor{black!45}{0.000} & \textcolor{black!45}{0.123} & 0.886 & \textcolor{black!45}{0.251} \\
BadNets 5 & \textbf{1.000} & 0.739 & 0.938 & 0.967 & \underline{0.999} & \underline{0.999} & 0.954 & 0.981 & 0.996 & 0.970 & \textcolor{black!45}{0.003} & 0.754 & \textcolor{black!45}{0.095} \\
BadNets 10 & \underline{1.000} & 0.985 & 0.988 & 0.927 & \textbf{1.000} & 1.000 & 0.966 & 0.982 & 0.994 & 0.992 & \textcolor{black!45}{0.000} & 0.665 & \textcolor{black!45}{0.013} \\
Blend 1 & 0.992 & \underline{0.998} & 0.869 & 0.728 & 0.638 & \textcolor{black!45}{0.468} & \textbf{0.999} & \textbf{0.999} & 0.943 & \textcolor{black!45}{0.249} & \textcolor{black!45}{0.241} & \textcolor{black!45}{0.282} & \textcolor{black!45}{0.001} \\
Blend 5 & \underline{0.999} & \textbf{0.999} & 0.983 & 0.867 & 0.741 & \textcolor{black!45}{0.480} & 0.999 & \underline{0.999} & 0.942 & \textcolor{black!45}{0.197} & \textcolor{black!45}{0.000} & \textcolor{black!45}{0.143} & \textcolor{black!45}{0.002} \\
Blend 10 & \underline{0.999} & \textbf{1.000} & 0.993 & 0.963 & 0.879 & 0.660 & \underline{0.999} & \textbf{1.000} & 0.972 & \textcolor{black!45}{0.001} & \textcolor{black!45}{0.000} & \textcolor{black!45}{0.267} & \textcolor{black!45}{0.001} \\
BPP 1 & \textbf{0.975} & \underline{0.974} & 0.781 & 0.936 & \textcolor{black!45}{0.104} & \textcolor{black!45}{0.105} & 0.969 & 0.973 & \textcolor{black!45}{0.454} & \textcolor{black!45}{0.124} & \textcolor{black!45}{0.003} & \textcolor{black!45}{0.422} & \textcolor{black!45}{0.399} \\
BPP 5 & \underline{0.986} & \textbf{0.987} & 0.872 & 0.912 & \textcolor{black!45}{0.009} & \textcolor{black!45}{0.010} & 0.966 & 0.986 & \textcolor{black!45}{0.233} & 0.650 & \textcolor{black!45}{0.190} & 0.776 & \textcolor{black!45}{0.381} \\
BPP 10 & \underline{0.991} & \textbf{0.996} & 0.815 & 0.952 & \textcolor{black!45}{0.020} & \textcolor{black!45}{0.007} & 0.832 & 0.832 & \textcolor{black!45}{0.341} & 0.900 & \textcolor{black!45}{0.248} & \textcolor{black!45}{0.221} & \textcolor{black!45}{0.187} \\
LF 1 & \textbf{0.935} & \underline{0.934} & 0.610 & \textcolor{black!45}{0.360} & \textcolor{black!45}{0.430} & \textcolor{black!45}{0.430} & 0.846 & 0.846 & \textcolor{black!45}{0.418} & 0.595 & \textcolor{black!45}{0.008} & \textcolor{black!45}{0.434} & \textcolor{black!45}{0.029} \\
LF 5 & \underline{0.988} & \textbf{0.989} & 0.925 & 0.748 & 0.565 & \textcolor{black!45}{0.311} & 0.967 & 0.967 & \textcolor{black!45}{0.222} & \textcolor{black!45}{0.048} & \textcolor{black!45}{0.001} & \textcolor{black!45}{0.157} & \textcolor{black!45}{0.008} \\
LF 10 & \textbf{0.977} & \underline{0.976} & 0.877 & 0.589 & \textcolor{black!45}{0.404} & 0.643 & 0.943 & 0.943 & \textcolor{black!45}{0.483} & \textcolor{black!45}{0.209} & \textcolor{black!45}{0.054} & \textcolor{black!45}{0.339} & \textcolor{black!45}{0.005} \\
WaNet 5 & \underline{0.928} & \textbf{0.938} & 0.521 & \textcolor{black!45}{0.218} & \textcolor{black!45}{0.117} & \textcolor{black!45}{0.125} & 0.771 & 0.771 & 0.889 & \textcolor{black!45}{0.394} & \textcolor{black!45}{0.006} & \textcolor{black!45}{0.226} & \textcolor{black!45}{0.037} \\
WaNet 10 & \textbf{0.970} & \underline{0.963} & 0.727 & \textcolor{black!45}{0.228} & \textcolor{black!45}{0.130} & \textcolor{black!45}{0.132} & 0.754 & 0.754 & 0.933 & 0.945 & \textcolor{black!45}{0.002} & 0.552 & \textcolor{black!45}{0.032} \\
\rowcolor{black!8}\emph{Tiny ImageNet mean} & \textbf{0.981} & \underline{0.951} & 0.796 & 0.738 & 0.502 & \textcolor{black!45}{0.455} & 0.925 & 0.929 & 0.701 & \textcolor{black!45}{0.448} & \textcolor{black!45}{0.063} & \textcolor{black!45}{0.437} & \textcolor{black!45}{0.103} \\
\rowcolor{black!8}\emph{all models} & \textbf{0.902} & 0.805 & \textcolor{black!45}{0.438} & 0.759 & 0.501 & \textcolor{black!45}{0.496} & \textcolor{black!45}{0.498} & \underline{0.869} & 0.602 & 0.662 & 0.715 & 0.788 & \textcolor{black!45}{0.166} \\
\bottomrule
\end{tabular}
\end{adjustbox}
\end{table}


\FloatBarrier

\subsection{Acceptance test}
\label{app:smoke}

Each port had to pass \cref{tab:detector-smoke} before it entered the comparison. It runs each detector on a benign, a BadNets and a Blend GTSRB model at 5\% poisoning with \SmokeValidationImages{} validation images (\SmokeValidationImagesExceptionSize{} for \SmokeValidationImagesExceptions{}), on the same splits and threshold rule as PSBD. Every detector reads chance on the benign model, so it raises no signal where there is no backdoor.

% generated by scripts/paper/tab_detector_smoke.py from docs/runs/2026-09-10-detector-smoke.md at 2026-09-29T16:25:20.399432+00:00, commit 19488b81358b06040ba96f361d5061b2981f1f25-dirty. Do not edit by hand.
\begin{table}[htbp]
\centering
\small
\caption{The acceptance test of each ported detector on three GTSRB models at 5\% poisoning, a benign one and two backdoored ones. The columns give AUROC at the headline quantile on each model, TPR at the smallest budget on the backdoored models, forward passes per input and seconds per input.}
\label{tab:detector-smoke}
\begin{adjustbox}{max width=\linewidth}
\begin{tabular}{lrrrrrrr}
\toprule
detector & passes & AUROC benign & AUROC BadNet & AUROC Blend & TPR@1\% BadNet & TPR@1\% Blend & s per input \\
\midrule
confidence & 1 & 0.503 & 0.738 & 0.876 & 0.000 & 0.034 & 0.001 \\
strip & 8 & 0.500 & 1.000 & 0.976 & 1.000 & 0.905 & 0.005 \\
scale\_up & 6 & 0.499 & 0.914 & 0.689 & 0.000 & 0.000 & 0.004 \\
scale\_up\_data\_limited & 6 & 0.499 & 0.909 & 0.678 & 0.000 & 0.000 & 0.002 \\
ibd\_psc & 6 & 0.500 & 0.481 & 0.970 & 0.000 & 0.724 & 0.002 \\
teco & 71 & 0.495 & 0.994 & 0.922 & 0.757 & 0.258 & 0.057 \\
cd\_l & 251 & 0.493 & 1.000 & 0.962 & 1.000 & 0.596 & 0.112 \\
beatrix & 1 & 0.510 & 0.453 & 0.819 & 0.000 & 0.000 & 0.001 \\
ted & 1 & 0.509 & 0.999 & 0.990 & 0.990 & 0.811 & 0.001 \\
sentinet & 202 & 0.499 & 0.107 & 0.379 & 0.056 & 0.051 & 0.057 \\
\bottomrule
\end{tabular}
\end{adjustbox}
\end{table}

\FloatBarrier
\clearpage


\section{The adaptive attacker per model}
\label{app:adaptive-per-model}

\Cref{tab:adaptive-attacker} gives every evasive model behind \cref{tab:adaptive}. For each model it lists the attack success after the evasion, the clean-accuracy cost, the AUROC of the probe the attacker trained against before and after, the mean AUROC of the two probes it never saw and the AUROC of the min-rank union of all three.

% generated by scripts/paper/tab_adaptive.py from results/adaptive_attacker_analysis.json, results/multi_probe_analysis.json at 2026-09-24T04:35:36.115870+00:00, commit 52e4ba1a7ac941365bc37d68e1e902dad4fb4f74-dirty. Do not edit by hand.
\begin{table}[htbp]
\centering
\small
\caption{The adaptive attacker on the models whose evasive version keeps attack success at or above 0.9.}
\label{tab:adaptive-attacker}
\begin{adjustbox}{max width=\linewidth}
\begin{tabular}{llllrrrrrr}
\toprule
arch & dataset & attack & rate & evade ASR & dCA & probed base & probed evade & transfer mean & union \\
\midrule
vit & CIFAR-100 & Adaptive-Blend & 0.05 & 0.948 & $-$0.026 & 0.927 & 0.034 & 0.936 & 0.944 \\
vit & CIFAR-100 & Adaptive-Blend & 0.1 & 0.925 & $-$0.045 & 0.951 & 0.039 & 0.921 & 0.927 \\
vit & CIFAR-100 & BadNets & 0.01 & 0.999 & $-$0.018 & 0.960 & 0.048 & 0.806 & 0.984 \\
vit & CIFAR-100 & BadNets & 0.05 & 1.000 & $-$0.026 & 0.977 & 0.004 & 0.785 & 0.994 \\
vit & CIFAR-100 & BadNets & 0.1 & 1.000 & $-$0.018 & 0.989 & 0.005 & 0.846 & 0.985 \\
vit & CIFAR-100 & Blend & 0.01 & 0.999 & $-$0.024 & 0.945 & 0.680 & 0.961 & 0.998 \\
vit & CIFAR-100 & Blend & 0.05 & 1.000 & $-$0.032 & 0.989 & 0.807 & 0.978 & 0.983 \\
vit & CIFAR-100 & Blend & 0.1 & 1.000 & $-$0.029 & 0.982 & 0.012 & 0.798 & 0.991 \\
vit & CIFAR-100 & BPP & 0.01 & 0.961 & $-$0.023 & 0.927 & 0.671 & 0.897 & 0.968 \\
vit & CIFAR-100 & BPP & 0.05 & 0.977 & $-$0.035 & 0.981 & 0.275 & 0.952 & 0.982 \\
vit & CIFAR-100 & BPP & 0.1 & 0.996 & $-$0.045 & 0.993 & 0.204 & 0.909 & 0.981 \\
vit & CIFAR-100 & LF & 0.01 & 0.938 & $-$0.042 & 0.944 & 0.410 & 0.838 & 0.852 \\
vit & CIFAR-100 & LF & 0.05 & 0.989 & $-$0.032 & 0.957 & 0.060 & 0.936 & 0.986 \\
vit & CIFAR-100 & LF & 0.1 & 0.996 & $-$0.026 & 0.994 & 0.014 & 0.938 & 0.967 \\
\bottomrule
\end{tabular}
\end{adjustbox}
\end{table}

\FloatBarrier
\clearpage


\section{Swin-S per model}
\label{supp:swin}

\Cref{tab:swin-attacks} gives every Swin-S model behind \cref{sec:swin}, and \cref{tab:swin-placements} every placement read on Swin with the number of models behind each.

% generated by scripts/paper/tab_swin.py from results/swin_*/psbd_metrics.json (90 implanted cells), checkpoints/swin_*/args.json at 2026-09-24T04:50:10.248016+00:00, commit 9300709d1f6e7a0ba702f153d77765353ffa1368-dirty. Do not edit by hand.
% generated by scripts/paper/tab_swin.py from results/swin_*/psbd_metrics.json (83 implanted cells), checkpoints/swin_*/args.json at 2026-09-24T06:22:48.655564+00:00, commit a6bece9dac4f610c3f14c6d149a3dd576da74d71-dirty. Do not edit by hand.
\begin{table}[htbp]
\centering
\small
\caption{Swin-S on CIFAR-10. Rows are attacks at their poison rate in percent. PSBD-TM masks tokens at the attention input and PSBD-RD, our adaptation of the original PSBD site, applies dropout after both residual additions. Both are read at the adaptive rule and the headline quantile, and the shaded row is the mean.}
\label{tab:swin-attacks}
\begin{adjustbox}{max width=\linewidth}
\begin{tabular}{llrrrr}
\toprule
attack & rate \% & PSBD-TM AUROC & TPR@10 & TPR@20 & PSBD-RD AUROC \\
\midrule
Adaptive-Blend & 5 & 0.496 & 0.03 & 0.10 & 0.787 \\
Adaptive-Blend & 10 & 0.992 & 0.99 & 0.99 & 0.981 \\
BadNets & 0.5 & 0.998 & 0.99 & 1.00 & 0.942 \\
BadNets & 1 & 0.993 & 1.00 & 1.00 & 0.839 \\
BadNets & 5 & 1.000 & 1.00 & 1.00 & 0.500 \\
BadNets & 10 & 1.000 & 1.00 & 1.00 & 0.996 \\
Blend & 0.5 & 0.991 & 0.97 & 1.00 & 1.000 \\
Blend & 1 & 0.990 & 0.97 & 0.99 & 0.999 \\
Blend & 5 & 0.997 & 1.00 & 1.00 & 0.994 \\
Blend & 10 & 1.000 & 1.00 & 1.00 & 1.000 \\
BPP & 0.5 & 0.985 & 0.96 & 0.98 & 0.989 \\
BPP & 1 & 0.800 & 0.38 & 0.60 & 0.995 \\
BPP & 5 & 0.997 & 0.99 & 0.99 & 0.999 \\
BPP & 10 & 0.998 & 1.00 & 1.00 & 0.998 \\
Label-Consistent & 10 & 0.965 & 0.95 & 0.96 & 0.530 \\
LF & 0.5 & 0.859 & 0.32 & 0.81 & 0.693 \\
LF & 1 & 0.996 & 0.99 & 0.99 & 0.989 \\
LF & 5 & 0.940 & 0.80 & 0.88 & 0.813 \\
LF & 10 & 0.982 & 0.93 & 0.96 & 0.953 \\
SIG & 5 & -- & -- & -- & 0.977 \\
SIG & 10 & 0.954 & 0.93 & 0.94 & 0.937 \\
TaCT & 1 & -- & -- & -- & 0.999 \\
TaCT & 5 & -- & -- & -- & 0.999 \\
WaNet & 5 & 0.980 & 0.85 & 0.98 & 0.961 \\
WaNet & 10 & 0.952 & 0.69 & 0.96 & 0.927 \\
\rowcolor{black!8}\emph{mean} &  & 0.948 & 0.85 & 0.92 & 0.912 \\
\bottomrule
\end{tabular}
\end{adjustbox}
\end{table}

% generated by scripts/paper/tab_swin.py from results/swin_*/psbd_metrics.json (83 implanted cells), checkpoints/swin_*/args.json at 2026-09-29T16:34:02.186194+00:00, commit 19488b81358b06040ba96f361d5061b2981f1f25-dirty. Do not edit by hand.
\begin{table}[htbp]
\centering
\small
\caption{Swin-S on CIFAR-100. Layout as in \cref{tab:swin-attacks}.}
\label{tab:swin-attacks-cifar-100}
\begin{adjustbox}{max width=\linewidth}
\begin{tabular}{llrrrr}
\toprule
attack & rate \% & PSBD-TM AUROC & TPR@10 & TPR@20 & PSBD-RD AUROC \\
\midrule
Adaptive-Blend & 5 & 0.921 & 0.91 & 0.91 & 0.920 \\
Adaptive-Blend & 10 & 0.970 & 0.97 & 0.97 & 0.973 \\
BadNets & 0.5 & 1.000 & 1.00 & 1.00 & 0.917 \\
BadNets & 1 & 1.000 & 1.00 & 1.00 & 0.643 \\
BadNets & 5 & 1.000 & 1.00 & 1.00 & 0.853 \\
BadNets & 10 & 1.000 & 1.00 & 1.00 & 0.927 \\
Blend & 0.5 & 0.986 & 0.96 & 0.99 & 0.993 \\
Blend & 1 & 0.999 & 1.00 & 1.00 & 0.992 \\
Blend & 5 & 1.000 & 1.00 & 1.00 & 0.998 \\
Blend & 10 & 0.999 & 1.00 & 1.00 & 0.999 \\
BPP & 0.5 & 0.985 & 0.96 & 0.98 & 0.977 \\
BPP & 1 & 0.997 & 0.99 & 1.00 & 0.994 \\
BPP & 5 & 0.998 & 1.00 & 1.00 & 0.995 \\
BPP & 10 & 0.998 & 0.99 & 1.00 & 0.996 \\
Label-Consistent & 1 & 0.999 & 1.00 & 1.00 & 0.716 \\
LF & 0.5 & 0.959 & 0.90 & 0.96 & 0.966 \\
LF & 1 & 0.945 & 0.85 & 0.95 & 0.874 \\
LF & 5 & 0.991 & 0.99 & 0.99 & 0.957 \\
LF & 10 & 0.996 & 0.99 & 1.00 & 0.986 \\
WaNet & 5 & 0.904 & 0.87 & 0.88 & 0.202 \\
WaNet & 10 & 0.927 & 0.89 & 0.93 & 0.803 \\
\rowcolor{black!8}\emph{mean} &  & 0.980 & 0.97 & 0.98 & 0.890 \\
\bottomrule
\end{tabular}
\end{adjustbox}
\end{table}

% generated by scripts/paper/tab_swin.py from configs/psbd_basis.json (panel rule retargeted to swin, 93 panel cells), results/swin_*/psbd_metrics.json (65 successful cells), checkpoints/swin_*/args.json at 2026-09-29T18:12:58.822484+00:00, commit d73d3277d007fd75b3a4166a7911b1e07a2761b7-dirty. Do not edit by hand.
\begin{table}[htbp]
\centering
\small
\caption{Swin-S on GTSRB. Layout as in \cref{tab:swin-attacks}.}
\label{tab:swin-attacks-gtsrb}
\begin{adjustbox}{max width=\linewidth}
\begin{tabular}{llrrrr}
\toprule
attack & rate \% & PSBD-TM AUROC & TPR@10 & TPR@20 & PSBD-RD AUROC \\
\midrule
Adaptive-Blend & 5 & 0.925 & 0.88 & 0.89 & 0.926 \\
Adaptive-Blend & 10 & 1.000 & 1.00 & 1.00 & 0.995 \\
BadNets & 1 & 1.000 & 1.00 & 1.00 & 0.729 \\
BadNets & 5 & 1.000 & 1.00 & 1.00 & 0.588 \\
BadNets & 10 & 1.000 & 1.00 & 1.00 & 0.989 \\
Blend & 1 & 0.978 & 0.92 & 0.94 & 0.999 \\
Blend & 5 & 1.000 & 1.00 & 1.00 & 1.000 \\
Blend & 10 & 0.987 & 0.95 & 0.97 & 1.000 \\
BPP & 1 & 0.979 & 0.93 & 0.99 & 0.991 \\
BPP & 5 & 0.997 & 0.99 & 1.00 & 1.000 \\
BPP & 10 & 0.998 & 0.99 & 1.00 & 0.999 \\
LF & 1 & 0.965 & 0.89 & 0.98 & 0.587 \\
LF & 5 & 0.980 & 0.94 & 0.97 & 0.670 \\
LF & 10 & 0.998 & 1.00 & 1.00 & 0.797 \\
WaNet & 5 & 0.885 & 0.87 & 0.87 & 0.301 \\
WaNet & 10 & 0.919 & 0.83 & 0.96 & 0.458 \\
\rowcolor{black!8}\emph{mean} &  & 0.976 & 0.95 & 0.97 & 0.814 \\
\bottomrule
\end{tabular}
\end{adjustbox}
\end{table}

% generated by scripts/paper/tab_swin.py from results/swin_*/psbd_metrics.json (83 implanted cells), checkpoints/swin_*/args.json at 2026-09-29T16:34:16.052244+00:00, commit 19488b81358b06040ba96f361d5061b2981f1f25-dirty. Do not edit by hand.
\begin{table}[htbp]
\centering
\small
\caption{Swin-S on Tiny ImageNet. Layout as in \cref{tab:swin-attacks}.}
\label{tab:swin-attacks-tiny-imagenet}
\begin{adjustbox}{max width=\linewidth}
\begin{tabular}{llrrrr}
\toprule
attack & rate \% & PSBD-TM AUROC & TPR@10 & TPR@20 & PSBD-RD AUROC \\
\midrule
Adaptive-Blend & 10 & 0.987 & 0.99 & 0.99 & 0.987 \\
BadNets & 1 & 0.998 & 1.00 & 1.00 & 0.651 \\
BadNets & 5 & 0.999 & 1.00 & 1.00 & 0.690 \\
BadNets & 10 & 1.000 & 1.00 & 1.00 & 0.848 \\
Blend & 1 & 1.000 & 1.00 & 1.00 & 0.998 \\
Blend & 5 & 1.000 & 1.00 & 1.00 & 0.998 \\
Blend & 10 & 1.000 & 1.00 & 1.00 & 0.999 \\
BPP & 1 & 0.997 & 1.00 & 1.00 & 0.996 \\
BPP & 5 & 0.999 & 1.00 & 1.00 & 0.999 \\
BPP & 10 & 1.000 & 1.00 & 1.00 & 0.999 \\
LF & 1 & 0.947 & 0.93 & 0.93 & 0.949 \\
LF & 5 & 0.995 & 1.00 & 1.00 & 0.993 \\
LF & 10 & 0.998 & 1.00 & 1.00 & 0.994 \\
WaNet & 5 & 0.991 & 0.99 & 0.99 & 0.834 \\
WaNet & 10 & 0.992 & 0.99 & 0.99 & 0.431 \\
\rowcolor{black!8}\emph{mean} &  & 0.994 & 0.99 & 0.99 & 0.891 \\
\bottomrule
\end{tabular}
\end{adjustbox}
\end{table}


\FloatBarrier
% generated by scripts/paper/tab_swin.py from configs/psbd_basis.json (panel rule retargeted to swin, 93 panel cells), results/swin_*/psbd_metrics.json (65 successful cells), checkpoints/swin_*/args.json at 2026-09-29T18:13:21.253463+00:00, commit d73d3277d007fd75b3a4166a7911b1e07a2761b7-dirty. Do not edit by hand.
\begin{table}[htbp]
\centering
\small
\caption{Mean AUROC on Swin-S at the headline quantile for every placement read on at least 5 backdoored models, at the matched and adaptive rules, over 65 models, 18 on CIFAR-10, 16 on CIFAR-100, 16 on GTSRB and 15 on Tiny ImageNet. Coverage is uneven across placements, so rows are not paired.}
\label{tab:swin-placements}
\begin{adjustbox}{max width=\linewidth}
\begin{tabular}{lrrrrr}
\toprule
placement & n match & AUROC matched & n adapt & AUROC adaptive & below chance \\
\midrule
token mask, attention output before the add & 65 & 0.804 & 65 & 0.850 & 8 \\
dropout, after both residual adds & 65 & 0.819 & 65 & 0.881 & 5 \\
dropout, stream after the attention add & 63 & 0.807 & 63 & 0.940 & 0 \\
token mask, stream after the attention add & 63 & 0.856 & 63 & 0.838 & 6 \\
dropout, embedding output & 63 & 0.840 & 63 & 0.878 & 5 \\
dropout, attention input after norm & 63 & 0.888 & 63 & 0.910 & 4 \\
dropout, attention input & 63 & 0.826 & 63 & 0.891 & 3 \\
channel mask, attention input & 63 & 0.839 & 63 & 0.893 & 3 \\
noise, attention input & 63 & 0.913 & 63 & 0.940 & 2 \\
token mask, attention input & 65 & 0.955 & 63 & 0.973 & 1 \\
noise, MLP input after norm & 63 & 0.798 & 63 & 0.727 & 16 \\
dropout, MLP input & 63 & 0.801 & 63 & 0.803 & 10 \\
token mask, MLP input & 63 & 0.728 & 63 & 0.828 & 7 \\
token mask, both sublayer inputs & 63 & 0.897 & 63 & 0.955 & 0 \\
scale up, input pixels & 63 & 0.760 & 63 & 0.814 & 9 \\
gain scale, MLP norm output & 63 & 0.887 & 63 & 0.828 & 10 \\
dropout, before both residual adds & 63 & 0.814 & 63 & 0.860 & 7 \\
dropout, before both residual adds, blocks 1 to 8 & 63 & 0.782 & 63 & 0.863 & 7 \\
dropout, before both residual adds, blocks 9 to 16 & 63 & 0.871 & 63 & 0.896 & 3 \\
token mask, attention input, blocks 1 to 8 & 63 & 0.745 & 53 & 0.802 & 9 \\
token mask, attention input, blocks 9 to 16 & 63 & 0.909 & 53 & 0.944 & 1 \\
dropout, before both residual adds, blocks 17 to 24 & 63 & 0.953 & 39 & 0.962 & 0 \\
dropout, MLP input after norm & 32 & 0.799 & 32 & 0.719 & 10 \\
dropout, MLP output before the add & 30 & 0.823 & 30 & 0.894 & 1 \\
token mask, embedding output & 26 & 0.658 & 26 & 0.684 & 7 \\
token mask, attention input, blocks 17 to 24 & 63 & 0.980 & 22 & 0.975 & 0 \\
token mask, stream after the MLP add & 16 & 0.863 & 16 & 0.846 & 1 \\
channel mask, attention input after norm & 16 & 0.743 & 16 & 0.809 & 2 \\
token mask, MLP input after norm & 16 & 0.730 & 16 & 0.726 & 5 \\
gain scale, attention norm output & 10 & 0.786 & 10 & 0.896 & 0 \\
\bottomrule
\end{tabular}
\end{adjustbox}
\end{table}

\FloatBarrier
\clearpage


\section{The mechanism in more detail}
\label{supp:mechanism}

\subsection{Per-layer profiles of the backdoor direction}
\label{app:layers}

\Cref{tab:tac-layers} gives, for each of the \TacLayersCheckpoints{} models in \cref{fig:tac-layers}, read on \TacLayersSamples{} paired images per model, the block at which the backdoor direction's norm relative to the clean activation peaks and that peak value. It also gives the onset, the first block at which the relative norm reaches half its final value, the final-layer linear centered kernel alignment (CKA) \cite{kornblith2019cka} between clean and triggered activations and the mean absolute trigger-activated change (TAC) per dimension at the final layer. A CKA of 1 means the trigger leaves the representation unchanged. The benign rows are the control. Their peak relative norm never exceeds \TacLayersBenignPeakMax{} and their CKA is at least \TacLayersBenignCkaMin{}, so the profiles of the backdoored rows are learned structure and not an artifact of the trigger's pixels. TaCT has a peak of only \TacLayersTactPeakMin{} to \TacLayersTactPeakMax{} and a CKA of \TacLayersTactCkaMin{} to \TacLayersTactCkaMax{}, because a source-specific attack changes the representation of its few source classes only and a direction averaged over all classes is small.

% generated by scripts/paper/mech_tac_layers.py from results/_experiments/tac_layers/tac_layers.json at 2026-09-29T18:01:43.340172+00:00, commit b2d32cf11708d5de965d3e13604c863a3ad9b493-dirty. Do not edit by hand.
\begin{table}[htbp]
\centering
\small
\caption{Per-layer profile of the backdoor direction on 15 ViT-B/16 models, 500 paired images each.}
\label{tab:tac-layers}
\begin{adjustbox}{max width=\linewidth}
\begin{tabular}{llrrrrr}
\toprule
Dataset & Model & Peak layer & Peak rel. norm & Onset layer & Final CKA & Final TAC \\
\midrule
GTSRB & BadNets 10\% & 12 & 1.64 & 7 & 0.57 & 1.05 \\
GTSRB & benign & 5 & 0.01 & 2 & 1.00 & 0.01 \\
GTSRB & Blend 10\% & 12 & 3.40 & 9 & 0.12 & 0.62 \\
GTSRB & BPP 10\% & 9 & 1.66 & 7 & 0.41 & 0.76 \\
GTSRB & LF 10\% & 12 & 1.68 & 7 & 0.48 & 0.86 \\
CIFAR-100 & BadNets 1\% & 12 & 0.79 & 8 & 0.44 & 0.57 \\
CIFAR-100 & BadNets 10\% & 10 & 0.92 & 8 & 0.39 & 0.57 \\
CIFAR-100 & benign & 10 & 0.12 & 6 & 0.99 & 0.06 \\
CIFAR-100 & Blend 1\% & 10 & 1.04 & 7 & 0.23 & 0.62 \\
CIFAR-100 & Blend 10\% & 9 & 1.70 & 5 & 0.16 & 0.71 \\
CIFAR-100 & BPP 10\% & 12 & 1.28 & 8 & 0.20 & 0.75 \\
CIFAR-100 & LF 10\% & 12 & 1.46 & 8 & 0.24 & 0.75 \\
Tiny ImageNet & BadNets 10\% & 12 & 0.99 & 9 & 0.41 & 0.55 \\
Tiny ImageNet & benign & 9 & 0.04 & 5 & 1.00 & 0.02 \\
Tiny ImageNet & Blend 10\% & 10 & 1.52 & 7 & 0.15 & 0.66 \\
\bottomrule
\end{tabular}
\end{adjustbox}
\end{table}

\FloatBarrier

\subsection{Comparison with Karayal\c{c}{\i}n et al.}
\label{sec:mechanism:picek}\label{app:picek}

Karayal\c{c}{\i}n et al.~\cite{karayalcin2026backdoor} studied backdoor directions in ViT-B/16 on the same three datasets with an overlapping set of attacks, and our models let us test three of their claims.

The first is that projecting the direction of one selected layer out of every weight matrix that writes to the residual stream, the patch embedding and every attention and MLP output projection, removes the backdoor. They report attack success falling from \cited{97.7} to \cited{6.7} percent on average, with Blend on CIFAR-100 as the one survivor. \Cref{tab:erasure} repeats their procedure with their layer selection rule and adds two controls, the same edit restricted to blocks 10 and 11 and the same edit with a random direction of the same norm. Over \ErasureCheckpoints{} checkpoints attack success falls from \ErasureAsrBeforeMean{} to \ErasureAsrAfterAllWritesMean{} and clean accuracy from \ErasureCaBeforeMean{} to \ErasureCaAfterAllWritesMean{}, and Blend on CIFAR-100 stays at \ErasureAsrAfterAllWritesMax{} as they report. Neither control moves the attack, with attack success at \ErasureAsrAfterBlocksOneZeroOneOneMean{} after the edit of blocks 10 and 11 and \ErasureAsrAfterRandomDirectionMean{} after the random direction. The direction is re-created by every block, so removing it means removing it everywhere, which is also why masking in every block beats masking in any third of the network (\cref{app:bands}).

\IfFileExists{tables/erasure.tex}{% generated by scripts/paper/tab_erasure.py from results/_experiments/whole_network_erasure/<folder>.json (10 checkpoints) at 2026-09-30T14:57:26.213851+00:00, commit 1fd38381a58031e98490ed24679d373b446725fb-dirty. Do not edit by hand.
\begin{table}[htbp]
\centering
\small
\caption{Attack success and clean accuracy after removing the backdoor direction from every weight that writes to the residual stream, on ViT-B/16 at 10\% poisoning, or at the rate named in the attack column where it differs.}
\label{tab:erasure}
\begin{adjustbox}{max width=\linewidth}
\begin{tabular}{llrrrrrrr}
\toprule
Dataset & Attack & ASR & CA & Layer & ASR, all writes & CA, all writes & ASR, blocks 10 and 11 & ASR, random direction \\
\midrule
CIFAR-10 & BadNets & 1.000 & 0.929 & 12 & 0.000 & 0.834 & 1.000 & 1.000 \\
CIFAR-10 & Blend & 1.000 & 0.937 & 12 & 0.000 & 0.864 & 1.000 & 1.000 \\
CIFAR-10 & BPP & 0.993 & 0.944 & 12 & 0.000 & 0.883 & 0.990 & 0.993 \\
CIFAR-10 & LF & 0.998 & 0.939 & 9 & 0.034 & 0.943 & 0.997 & 0.998 \\
CIFAR-10 & WaNet & 0.876 & 0.933 & 7 & 0.009 & 0.931 & 0.855 & 0.884 \\
CIFAR-100 & BadNets & 1.000 & 0.807 & 11 & 0.009 & 0.810 & 1.000 & 1.000 \\
CIFAR-100 & Blend (5\%) & 1.000 & 0.833 & 9 & 0.991 & 0.833 & 1.000 & 1.000 \\
CIFAR-100 & Blend & 1.000 & 0.820 & 7 & 1.000 & 0.819 & 1.000 & 1.000 \\
CIFAR-100 & BPP & 0.988 & 0.824 & 12 & 0.000 & 0.820 & 0.985 & 0.988 \\
CIFAR-100 & WaNet & 0.801 & 0.827 & 7 & 0.000 & 0.834 & 0.773 & 0.808 \\
\bottomrule
\end{tabular}
\end{adjustbox}
\end{table}
}{}
\FloatBarrier

The second is that distributed triggers reach the class token earlier than patch triggers. In \cref{fig:tac-layers} the backdoor direction's onset comes earlier for Blend, BPP and LF than for BadNets on CIFAR-100 and Tiny ImageNet, and the order reverses on GTSRB, which they did not run.

The third is a data-free detector that scores the classifier head's rows against the early layers' output projections and flags a model when one class stands out by more than 3 standard deviations, reported to flag WaNet and BPP reliably and patch triggers never. We ran it with their grid over the number of layers and the floor of the standard deviation on \WeightDetectorBackdoored{} backdoored and \WeightDetectorBenign{} benign models, and \cref{tab:weight-detector} gives every model. It names the correct target on all \WeightDetectorWanetModels{} WaNet models, with a best $Z$ of \WeightDetectorWanetBestZMin{} to \WeightDetectorWanetBestZMax{}, on \WeightDetectorBppNamed{} of \WeightDetectorBppModels{} BPP models and on the \WeightDetectorLfModels{} LF model. It never names the target of a BadNets, Blend or TaCT model. On the benign models \WeightDetectorBenignFlagMin{} to \WeightDetectorBenignFlagMax{} percent of the grid flags, every flag a false alarm. The backdoor is a rotated direction rather than a coordinate pattern (\cref{tab:direction-ablation}), so a score built from coordinates of the head and the weights sees only its projection.

% generated by scripts/paper/app_weight_detector.py from results/_experiments/weight_detector/summary.json at 2026-09-24T06:05:56.925219+00:00, commit a6bece9dac4f610c3f14c6d149a3dd576da74d71-dirty. Do not edit by hand.
\begin{table}[htbp]
\centering
\small
\caption{The data-free weight detector of Karayalcin et al. over their grid on our models. Named is whether some grid cell flags the model and names its true target, best Z the largest Z among those cells and flag share the share of the grid that flags, which on a benign model is the false-alarm rate.}
\label{tab:weight-detector}
\begin{adjustbox}{max width=\linewidth}
\begin{tabular}{llrlrr}
\toprule
dataset & attack & rate \% & named & best Z & flag share \\
\midrule
CIFAR-100 & BadNets & 10 & no & -- & 0.04 \\
CIFAR-100 & benign & -- & -- & -- & 0.15 \\
CIFAR-100 & Blend & 10 & no & -- & 0.10 \\
CIFAR-100 & BPP & 1 & yes & 9.9 & 0.08 \\
CIFAR-100 & BPP & 5 & no & -- & 0.06 \\
CIFAR-100 & BPP & 10 & no & -- & 0.06 \\
CIFAR-100 & LF & 10 & yes & 10.0 & 0.05 \\
CIFAR-100 & WaNet & 5 & yes & 139.3 & 0.27 \\
CIFAR-100 & WaNet & 10 & yes & 69.6 & 0.16 \\
CIFAR-10 & BadNets & 10 & no & -- & 0.13 \\
CIFAR-10 & benign & -- & -- & -- & 0.03 \\
CIFAR-10 & BPP & 10 & yes & 71.4 & 0.07 \\
CIFAR-10 & WaNet & 10 & yes & 40.0 & 0.07 \\
Tiny ImageNet & BadNets & 10 & no & -- & 0.13 \\
Tiny ImageNet & benign & -- & -- & -- & 0.11 \\
Tiny ImageNet & BPP & 10 & no & -- & 0.07 \\
Tiny ImageNet & WaNet & 10 & yes & 34.6 & 0.15 \\
\bottomrule
\end{tabular}
\end{adjustbox}
\end{table}

\FloatBarrier
